\documentclass[10pt,twocolumn]{article}
\usepackage{amsmath,amssymb,amsthm}
\usepackage[margin=0.85in]{geometry}
\usepackage{booktabs}
\usepackage{hyperref}
\hypersetup{colorlinks=true,linkcolor=blue,citecolor=blue,urlcolor=blue}

\theoremstyle{plain}
\newtheorem{theorem}{Theorem}
\newtheorem{proposition}[theorem]{Proposition}
\newtheorem{lemma}[theorem]{Lemma}
\newtheorem{corollary}[theorem]{Corollary}
\theoremstyle{definition}
\newtheorem{reported}[theorem]{Reported computational result}
\newtheorem{conditional}[theorem]{Conditional corollary}
\newtheorem{conjecture}[theorem]{Conjecture}
\theoremstyle{remark}
\newtheorem{remark}[theorem]{Remark}

% Every displayed result that a machine check establishes names it.
\newcommand{\chk}[1]{\par\smallskip\noindent{\footnotesize\textsf{check:}
  \texttt{#1}}\par\smallskip}

\newcommand{\CF}{C_F}
\newcommand{\dg}{^{\dagger}}

\begin{document}

\title{\bfseries Nested-quotient temporal histories and homological flat bands
in strong-coupling $SU(N)$ Hamiltonian lattice gauge theory\\[2mm]
\large Final edition, with machine-checked provenance}
\author{Alexander Smith\\ \small Independent Researcher}
\date{29 August 2026}
\maketitle

\begin{abstract}
\noindent
We study the Kogut--Susskind Hamiltonian for pure $SU(N)$ gauge theory on a
periodic cubic spatial lattice at strong coupling. After selecting the
complete gauge-invariant charge-odd one-plaquette eigenspace, fixed-order
contributions are organised as ordered magnetic histories on the Haar--Gram
quotient of reachable Wilson networks. What is exact and all-rank is one matrix
element: the connected second-order coupling between two plaquettes sharing an
oriented link is $t_Nu^2$ times their incidence overlap, with
\[
  t_N=\frac{2N(N^2-4)}{(N^2-1)(2N^2-1)(4N^2-9)}>0,\qquad N\ge 3 .
\]
Under two stated hypotheses --- that the retained shell is exactly the
charge-odd one-plaquette span, and that no connected non-adjacent history
carries momentum dependence --- those elements assemble into the oriented
edge--face Gram operator $t_Nu^2B(k)B(k)\dg$ with $B=\partial_2\dg$, whose
projected spectrum is $\{0,\,t_Nu^2q,\,t_Nu^2q\}$ up to a common scalar,
$q=4\sum_j\sin^2(k_j/2)$. Neither hypothesis is checked here and both are named
where they are used. The flat carrier is $\ker\partial_2$, of exact dimension
$L^3+2$, and that is unconditional.

Relative to the 28 August edition this one adds five things, each carrying its
own machine check. The channel decomposition of $t_N$, previously a
representation-theoretic ledger derived on a single plaquette pair, is recovered on a
genuine face-sharing two-cube $SU(3)$ cluster --- an exact-Casimir $B=6$ basis
of $1{,}590{,}462$ states, which the delivery reaches by path enumeration
without ever assembling the magnetic matrix --- and the
two decompositions agree \emph{channel by channel} rather than only in the sum
(Reported result~\ref{rep:census}); on that space each of the six channels is
separately proportional to the same incidence Gram, on all $56$ adjacent pairs,
so the separation of colour from geometry is measured and not assumed, and the
connected diagonal --- eleven numbers in both deliveries --- is constant on the
prism's three symmetry orbits, with the whole $B=4\to B=6$ difference confined
to the eight faces that carry connected transport
(Proposition~\ref{prop:orbit}).
Separately, and by arithmetic on the weight
ledger alone, a cutoff keeping only
the singlet and $\Lambda^2F$ routes projects the hopping to $-1/12$ ---
opposite in sign to $t_3$ --- and omits exactly $14/153$, the two summing to
$t_3$ (Corollary~\ref{cor:cutoff}); the $B=4$ truncation of the two-cube space
delivers the first of those two numbers outright, on the same geometry. The charge-even sector acquires a closed form:
with $\hat a_m=2+2\cos k_m$ and $p=\sum_m\hat a_m=12-q$, its Bloch spectrum is
the cubic $\mu(\mu-p)^2=4\hat a_1\hat a_2\hat a_3$, so one zone function runs
both sectors (Theorem~\ref{thm:evencubic}), and its
$\Gamma$-point splitting pins the charge-even hopping $t_{r,+}$ --- a
separation the charge-odd rest frame has no analogue for, since its
$\Gamma$ spectrum has only one level (Proposition~\ref{prop:evengamma}). The
homological reading of the flat direction is shown to be the same one on a
second complex: the six faces of a cube are a closed surface, its
$\ker\partial_2$ is the fundamental class, and the one-cube charge-odd shell
splits as $A_1^{--}\oplus T_1^{+-}\oplus E^{--}$ with the flat level sitting at
the scalar for \emph{every} value of the hopping
(Proposition~\ref{prop:cube}). Finally the open
fourth-order off-axis coefficient is given a geography rather than a verdict:
an exact obstruction certificate, an explicit second witness, and the
measurement that no $\Gamma$-anchoring error can move it at all
(Section~\ref{sec:fourth}).

The two-cube channel census, and the $SU(3)$ third- and fourth-order
statements, remain conditional on supplied ledgers --- the first on six
rational reconstructions of finite-precision contractions, the others on
exact-arithmetic ledgers --- whose raw payloads are not reproduced here, and
all are labelled as such throughout. Every displayed result carries the name of the
machine check that establishes it. All statements concern a formal
finite-order projected sector: this is not a continuum glueball or mass-gap
theorem.
\end{abstract}

\section{Introduction}

Hamiltonian strong-coupling expansions organise lattice gauge dynamics around
the electric vacuum and its flux excitations~\cite{KS1975,KSS1976}. In three
spatial dimensions a fundamental plaquette excitation carries one of three face
orientations, and at second order neighbouring faces communicate through a
shared link, so the projected problem resembles a three-orbital tight-binding
model. The signs of those matrix elements are fixed by the orientation of the
plaquette boundary and by charge conjugation, and the resulting interference is
the central object here.

Flat bands generated by incidence or Gram matrices have a long history: local
topology, line-graph constructions, compact localised states, and the need for
noncontractible states on periodic systems are all established, and closely
related closed-surface structures occur in two-form spin liquids. A recent
classification of exactly flat bands protected by local graph
topology~\cite{Hazra2026} reaches extensive face-graph degeneracy from an
\emph{unoriented} face--vertex incidence, where rank--nullity alone gives
$|F|-|V|$ flat bands. That is a different mechanism from the one here, and the
difference is measurable: on the signed face--edge operator of
Section~\ref{sec:carrier} the corresponding index is exactly $0$ and bounds
nothing, while both kernels equal $L^3+2$ --- homology, not counting. It is
recorded as a distinction, not an identification. We do not claim that
incidence flatness, cellular compact localised states, or homological counting
are new in isolation.
\chk{the discrete index theorem gives 0 on the signed face-edge operator, bounding nothing}

The content is dynamical. Starting from non-Abelian $SU(N)$ Hamiltonian
lattice gauge theory, exact strong-coupling matrix elements generate an
oriented edge--face Gram operator in the charge-odd one-plaquette sector.
Representation theory fixes its coefficient and the cubical chain complex
fixes its kernel and spectrum. The separation is exact,
\[
  \begin{aligned}
    \text{colour dynamics}&\longrightarrow t_N,\\
    \text{cellular geometry}&\longrightarrow BB\dg .
  \end{aligned}
\]

\subsection*{What is proved, and what is conditional}

Nine statements are self-contained here.

\begin{enumerate}\itemsep2pt
\item For every $N\ge3$ the shared-link channel weights are $d_\rho/N^2$,
  derived from the order-two Weingarten values
  (Theorem~\ref{thm:weingarten}). This was previously an isotropy
  \emph{hypothesis}, and everything in Section~\ref{sec:second} rests on it.
\item For every $N\ge3$ the connected second-order coupling between two
  plaquettes sharing one link is $t_Nu^2$ times their incidence overlap
  (Theorem~\ref{thm:allrank}), and the channel list it runs through is complete
  (Lemma~\ref{lem:channels}). The assembly of those elements into the torus
  operator is \emph{not} here; it is item~\ref{item:assembly} below.
\item The Bloch spectrum of $BB\dg$ is $\{0,q,q\}$ (Theorem~\ref{thm:fact}),
  and on the periodic $L^3$ complex the zero-mode count is $L^3+2$
  (Theorem~\ref{thm:carrier}). The two are the same statement, by a sum rule
  over the momentum grid (Theorem~\ref{thm:bridge}) --- and not merely in
  agreement: the $\Gamma$ block is the harmonic subspace, so its $3$ is
  $b_2(T^3)$, and the wrapping sheets averaged over their translates
  \eqref{eq:harm} are harmonic representatives of the three classes.
\item\label{item:cutoff} The truncation that keeps only the singlet and $\Lambda^2F$ routes
  projects the hopping to $w_{\Lambda^2F}-w_{\mathbf 1}=-1/12$, and what it
  omits is $w_{\mathrm{Sym}^2F}-w_{\mathrm{Adj}}=14/153$; the two sum to
  $t_3$ (Corollary~\ref{cor:cutoff}). This is arithmetic on
  \eqref{eq:w} and needs no input beyond Theorem~\ref{thm:weingarten}.
\item The charge-even Bloch spectrum is the cubic
  $\mu(\mu-p)^2=4\hat a_1\hat a_2\hat a_3$ with $p+q=12$
  (Theorem~\ref{thm:evencubic}); its range is exactly $[-4,12]$, the top
  attained only at $\Gamma$ and the floor on the three planes $k_j=\pi$, triple
  only at $R$ (Proposition~\ref{prop:evenrange}); and its $\Gamma$-point
  splitting is $12t_{r,+}$ exactly (Proposition~\ref{prop:evengamma}). The
  identity is self-contained; the two numbers it is evaluated at inherit the
  status of $t_{r,+}$, and the $r=3$ one is conditional --- see
  item~\ref{item:domino}. That the charge-odd rest frame is blind to the hopping
  and the charge-even one is not follows from Schur's lemma alone: the
  charge-odd triple carries the irreducible $T_1$ and the charge-even one
  $A_1\oplus E$, so the two $\Gamma$ blocks have one level and two, at every
  order and independently of the incidence ansatz.
\item The vacuum-mediated route $p\to|0\rangle\to p'$ connects any pair of
  plaquettes and is worth $1/\CF$; it vanishes in the charge-odd sector because
  $V$ is charge-even, which is what makes Theorem~\ref{thm:allrank}'s
  shared-link restriction true here and false in general, and it gives the
  all-rank charge-even hopping $\ell_N$ in closed form
  (Proposition~\ref{prop:vacuum}).
\item A correction of the boundary-factorised form
  $a_rI+b_rS+BM_rB\dg$ shifts the carrier without dispersing it, for
  \emph{every} $M_r$ (Proposition~\ref{prop:rigid}). The protection is a
  property of the boundary operator, not of the dynamics --- and the
  proposition says nothing about a correction that does not factor through
  links, which is why the fourth order has to be argued.
\item On the closed cube surface the charge-odd shell is
  $A_1^{--}\oplus T_1^{+-}\oplus E^{--}$, with $G$-eigenvalues $0,4,6$; the
  labels follow from the rotation action rather than being assigned, the
  $A_1^{--}$ level is the fundamental class of $S^2$, and it sits at the scalar
  for every hopping coefficient (Proposition~\ref{prop:cube}). This is integer
  linear algebra and takes no input from any delivery.
\item The retained $\Gamma$/axis corpus cannot identify the disputed
  fourth-order off-axis coefficient (Theorem~\ref{thm:obstruction}), and an
  explicit second witness exhibits the non-identifiability by construction
  (Proposition~\ref{prop:witness}).
\end{enumerate}

Two families of statements are \emph{conditional} on supplied exact-arithmetic
derivation ledgers. Their raw class contractions, symbolic history
certificates and microscopic rank-sweep artifacts are not included here, and
repeated assertions in synthesis documents are not treated as independent
reproductions.

\begin{enumerate}\setcounter{enumi}{9}\itemsep2pt
\item\label{item:assembly} On two hypotheses that no check covers --- that the
  range of $P_-$ is exactly the charge-odd one-plaquette span, and that no
  connected non-adjacent second-order history carries momentum dependence ---
  the adjacent elements of item~2 assemble into $t_Nu^2BB\dg$, with the band
  structure that follows from it (Conditional corollary~\ref{cor:assembly}).
  Everything downstream that quotes a bandwidth or a band edge at order $u^2$
  inherits those two hypotheses.
\item A supplied two-cube delivery reports six link-irrep channel
  coefficients at $N=3$, as rational reconstructions of finite-precision
  contractions. Conditional on those six numbers, they are the four
  fusion-channel weights individually --- channel by channel and not merely in
  the sum (Reported result~\ref{rep:census}) --- and the $B=4$ truncation of
  that space retains exactly the adjacent shared-link channels the published
  link cutoff of item~\ref{item:cutoff} retains, so the two land on the same
  projected hopping --- the same channels, not the same scheme. Conditional on the delivered diagonals, the symmetry
  analysis of Proposition~\ref{prop:orbit} applies to them: the orbit structure
  itself is exact and takes no delivered input.
  The $B=6$ two-cube build was not reproduced here.
\item\label{item:domino} For $SU(3)$, supplied ledgers report that all
  candidate third-order lifter classes vanish and report the scalar and shared-link coefficients of
  Section~\ref{sec:su3}. The census is three lifter classes evaluated over
  $32$ five-trace numerator patterns each, $96$ contractions in total, and
  nothing here checks any of them. If it is exhaustive and the linked ledger
  complete, the effective Hamiltonian stays in the incidence algebra through
  order $u^3$.
\item Supplied ledgers report the all-rank cube formula
  $\alpha_N=640/[N(N^2-1)^3]$ and its equality with the complete axial
  coefficient for $N=3,\dots,12$. Conditional on those calculations, the
  $SU(3)$ carrier acquires nonzero axial dispersion at order $u^4$. The
  complete off-axis operator remains open --- and by
  Theorem~\ref{thm:obstruction} it cannot be closed by any $\Gamma$ or axial
  datum.
\end{enumerate}

These statements concern the effective Hamiltonian projected to the degenerate
one-plaquette manifold. Virtual multi-plaquette states are retained in the
resolvents, but the external spectral problem is not the full gauge-theory
Hilbert space. At $k=0$ the three incidence branches touch and the
finite-volume separation collapses as $L^{-2}$. Nothing below proves
convergence of the strong-coupling series, persistence under unrestricted
higher-order operators, identification with the lightest physical glueball, or
a continuum mass gap.

\section{Hamiltonian and projected sector}
\label{sec:sector}

Let $\Lambda_L=(\mathbb{Z}/L\mathbb{Z})^3$, $L\ge3$, be a cubic spatial
lattice with periodic boundary conditions. In units of the electric coupling
the Kogut--Susskind Hamiltonian is
\begin{equation}
  H_\beta=\tfrac12\sum_\ell C_2(\ell)
  +\beta_N\sum_p\Bigl(1-\tfrac1N\operatorname{Re}\operatorname{Tr}U_p\Bigr).
\end{equation}
Dropping the additive constant,
\begin{equation}\label{eq:H}
  H=H_E-uV,\qquad V=\sum_p(\chi_p+\bar\chi_p),\qquad u=\frac{\beta_N}{2N},
\end{equation}
with $\chi_p=\operatorname{Tr}U_p$. For $SU(3)$, $u=\beta_3/6$. The
definition $u=\beta_N/(2N)$, not the $SU(3)$-specific equality, is the
all-rank convention; the archived line $Y=2\beta/3=4u$ is a definition-label
erratum and the printed coefficients are already in $u$. Reading it as a
coordinate would demand dividing the order-$r$ coefficient by $4^r$, which
misses the printed tables by exactly $16$ at $u^2$ and $64$ at $u^3$.
\chk{the printed towers are canonical-u: 4*Delta(3u/2) reproduces them verbatim}
\chk{the 4**r rescaling breaks the bridge: order 2 off by 16, order 3 by 64}

All quoted excitation energies use the linked, vacuum-subtracted operator
$H_{\mathrm{exc}}(u)=H(u)-E_{\mathrm{vac}}(u)I$. This changes only the scalar
ledger; incidence hopping and carrier shape are unaffected.

The electric vacuum $|0\rangle$ carries the trivial representation on every
link. For $N\ge3$ the normalised one-plaquette states of definite charge
conjugation are
\begin{equation}\label{eq:states}
  |p,\pm\rangle=\frac{\chi_p\pm\bar\chi_p}{\sqrt2}\,|0\rangle ,
\end{equation}
of unit norm by character orthogonality. Each of the four boundary links
carries the fundamental or antifundamental representation, so
\begin{equation}\label{eq:EF}
  E_{F,N}=4\Bigl(\tfrac12 \CF\Bigr)=2\CF=\frac{N^2-1}{N},
  \qquad \CF=\frac{N^2-1}{2N}.
\end{equation}
This external energy is the input to the resolvent denominators of
Section~\ref{sec:second}, where it is checked together with them.
For $SU(2)$ complex conjugation is an inner gauge transformation and the
charge-odd state vanishes; this is why the construction begins at $N=3$. The
same fact appears again in Section~\ref{sec:second}: $\Lambda^2F$ is the
singlet exactly at $N=2$, which is where $t_N$ vanishes. The charge-even
hopping does \emph{not} vanish there, which locates the $N^2-4$ zero in the
channel difference rather than in the incidence.
\chk{at N = 2 the C-odd hopping vanishes and the C-even one does not}

We work inside the charge-odd fundamental Wilson-loop cyclic sector with
trivial electric centre flux around each torus cycle. Let $P_-$ be the full
$E_{F,N}$ spectral projector there. For $L\ge3$ its range is the charge-odd
one-plaquette span, since a reduced contractible fundamental loop at that
energy has four distinct links and is an elementary plaquette. This sector
qualification also excludes the accidentally degenerate straight wrapping loop
at $L=4$. With $Q=1_--P_-$ the reduced resolvent crosses no omitted
$E_{F,N}$ eigendirection, and degenerate perturbation theory gives
\begin{equation}\label{eq:heff}
  H_{\mathrm{eff},-}=P_-H_EP_-+uH_1+u^2H_2+u^3H_3+\cdots
\end{equation}
with every $H_r$ linked, vacuum-subtracted, and carrying the appropriate
$Q$-space resolvents.

\section{Oriented incidence and the exact carrier}
\label{sec:carrier}

The periodic cubical complex has the cellular chain sequence
$C_3\xrightarrow{\partial_3}C_2\xrightarrow{\partial_2}C_1$ with
$\partial_2\partial_3=0$. One-plaquette amplitudes live in $C_2$ and their
signed leakage into links is $\partial_2c$: a cellular closed-surface
condition, not the microscopic vertex Gauss law, which the Wilson-loop basis
already satisfies through $\partial_1\partial_2=0$.

Writing $d_j=e^{ik_j}-1$ and $q(k)=\sum_j|d_j|^2=4\sum_j\sin^2(k_j/2)$, the
face-to-link Bloch incidence matrix is
\begin{equation}\label{eq:B}
  B(k)=\partial_2(k)\dg=
  \begin{pmatrix}
    d_2 & -d_1 & 0\\
    d_3 & 0 & -d_1\\
    0 & d_3 & -d_2
  \end{pmatrix}.
\end{equation}

\begin{theorem}[Incidence factorisation]\label{thm:fact}
With $w(k)=(\bar d_3,\,-\bar d_2,\,\bar d_1)^{\mathsf T}$,
\[
  BB\dg=qI-ww\dg,\qquad B\dg w=0,\qquad \|w\|^2=q,
\]
and therefore $\operatorname{spec}(BB\dg)=\{0,q,q\}$. With
$S=BB\dg-4I$ the signed shared-link adjacency,
$\operatorname{spec}S(k)=\{-4,\,q-4,\,q-4\}$.
\end{theorem}

\begin{proof}
Multiplying \eqref{eq:B} by its adjoint gives the first identity entry by
entry. Since $\|w\|^2=q$, the rank-one operator $qI-ww\dg$ annihilates $w$ and
acts as $qI$ on its orthogonal complement. Subtracting $4I$ gives the
adjacency spectrum.
\end{proof}
\chk{B B\^{}dagger = q I - d conj(d)\^{}T for the curl incidence}

For $k\neq0$ the carrier is one-dimensional and represented by $w(k)$. At
$\Gamma$, $B=0$ and all three branches meet; the normalised Bloch eigenvector
has no direction-independent limit there, which is the familiar
singular-flat-band mechanism behind the incompleteness of a purely compact
localised basis~\cite{BWB2008}.

\begin{theorem}[Finite-torus carrier]\label{thm:carrier}
For the periodic cubic cellulation of $T^3$,
$\ker\partial_2=\operatorname{im}\partial_3\oplus H_2$ with
$\dim\operatorname{im}\partial_3=L^3-1$ and $\dim H_2=3$, so
$\dim\ker\partial_2=L^3+2$.
\end{theorem}

\begin{proof}
There are $L^3$ three-cells and no four-cells, and $b_3(T^3)=1$, so
rank$\,\partial_3=L^3-1$ by rank--nullity. Discrete Hodge decomposition gives
$\ker\partial_2=B_2\oplus H_2$ with $\dim H_2=b_2(T^3)=3$.
\end{proof}

The count is verified over $\mathbb{Z}$ rather than re-derived from the
formula: the complex is built from the boundary formulas of
Appendix~\ref{app:chain}; the exhibited cycles --- every elementary cube
boundary together with three wrapping sheets --- bound the kernel from below,
and rank over $\mathbb{F}_p$ bounds it from above, since rank can only drop
under reduction. The two bounds meet at $L=1,\dots,5$.
\chk{dim Z\_2 = L\^{}3 + 2 by rank, not by re-arranging the formula}
\chk{cube boundaries and three wrapping sheets SPAN Z\_2}

\begin{remark}
$L^3+2$ is the nullity of the down boundary Laplacian, not a Betti number; the
Betti contribution is the final $3$. The restriction $L\ge3$ belongs to the
twelve-neighbour Bloch adjacency, where $x+e$ and $x-e$ coincide at $L=2$. The
chain-level count has no such restriction and holds at $L=1,2$ as well.
\chk{the L\^{}3+2 count is chain-level, not the Bloch convention}
\end{remark}

The following sum rule joins Theorems~\ref{thm:fact} and~\ref{thm:carrier},
which are otherwise two unrelated statements: a $3\times3$ spectrum and an
$(L^3+2)$-dimensional space.

\begin{theorem}[Bloch--chain bridge]\label{thm:bridge}
On the periodic $L^3$ complex,
\[
  \sum_{k}\bigl(3-\operatorname{rank}B(k)\bigr)=L^3+2 ,
\]
the sum running over the $L^3$ allowed momenta, with nullity profile
$3$ at $\Gamma$ exactly once and $1$ at each of the other $L^3-1$ momenta.
\end{theorem}

\begin{proof}
$B(0)=0$ contributes $3$. For $k\neq0$ some $d_j\neq0$, and
Theorem~\ref{thm:fact} gives $\operatorname{rank}B(k)=2$, contributing $1$
each. Summing gives $3+(L^3-1)=L^3+2$.
\end{proof}
\chk{the Bloch and chain routes to the carrier agree}

\begin{remark}
Earlier editions of this paper read the agreement as a coincidence worth
checking --- ``the $3$ is the triple degeneracy of $B(0)=0$, not $b_2(T^3)$''.
That was wrong, and the correction is the better statement. Under the discrete
Fourier transform $\partial_2$ block-diagonalises and
$\ker\partial_2=\bigoplus_k\ker\partial_2(k)$; at $k\neq0$ the kernel block is
one-dimensional and is exactly $\operatorname{im}\partial_3(k)$, while
$\partial_3(0)=0$. So all of $\operatorname{im}\partial_3$ sits at nonzero
momenta and the quotient is supported entirely at $\Gamma$: the $\Gamma$ block
\emph{is} $H_2$, and the $3$ \emph{is} $b_2(T^3)$. The two routes do not agree
by accident; they are the same decomposition read in two bases.
\chk{the sheets average to harmonic representatives, and the Gamma block IS b\_2}
\end{remark}

\subsection{The wrapping sheets are cycles, not harmonic --- and their
harmonic representatives}

Theorem~\ref{thm:carrier} exhibits its three $H_2$ generators explicitly as
wrapping sheets $s_{(ij)}=\sum_{x:\,x_m=0}[x;i,j]$. These satisfy
$\partial_2 s=0$ exactly, so they are cycles and they span the quotient
$\ker\partial_2/\operatorname{im}\partial_3$. They are \emph{not} harmonic:
$\partial_3\dg s\neq0$, and
\begin{equation}
  \frac{\langle s,\;\partial_3\partial_3\dg\,s\rangle}{\|s\|^2}=2
  \qquad\text{exactly, for every }L\ge2 .
\end{equation}
\chk{FINDING: the wrapping sheets are cycles but NOT harmonic}

The distinction matters because $H_2$ is used in two senses: the quotient in
Theorem~\ref{thm:carrier}, which the sheets span, and the harmonic subspace
$\ker\partial_2\cap\ker\partial_3\dg$, which they do not lie in.

Harmonic representatives do exist, and they are one line away. Averaging a
sheet over its $L$ translates in the complementary direction,
\begin{equation}\label{eq:harm}
  h_{(ij)}=\frac1L\sum_{r=0}^{L-1}s_{(ij)}(r),
  \qquad
  \partial_2h_{(ij)}=0,\quad \partial_3\dg h_{(ij)}=0 ,
\end{equation}
gives three harmonic chains spanning
$\ker\partial_2\cap\ker\partial_3\dg$, which is exactly three-dimensional;
and $s_{(ij)}-h_{(ij)}\in\operatorname{im}\partial_3$, so a sheet and its
average are the same homology class. The averaging is the $k=0$ Fourier
projection, and what it removes is precisely the non-harmonic part measured
above. Nothing about the finding is withdrawn --- the generators
Theorem~\ref{thm:carrier} exhibits are still not harmonic --- but a real-space
statement phrased for the harmonic subspace now has representatives to point
at, and Proposition~\ref{prop:rigid} covers both, since $B\dg w=0$ holds on
all of $\ker\partial_2$.
\chk{the sheets average to harmonic representatives, and the Gamma block IS b\_2}

\section{Exact all-rank dynamics at second order}
\label{sec:second}

Two plaquettes sharing a link leave six nonshared fundamental links after the
external contractions. On the shared link the fusion channels depend on the
relative induced orientation,
\begin{equation}\label{eq:fusion}
  F\otimes\bar F=\mathbf{1}\oplus\mathrm{Adj},
  \qquad
  F\otimes F=\Lambda^2F\oplus\mathrm{Sym}^2F ,
\end{equation}
with dimensions and quadratic Casimirs
\begin{equation}\label{eq:table}
\begin{aligned}
  d_{\mathbf 1}&=1, & C_{\mathbf 1}&=0,\\
  d_{\mathrm{Adj}}&=N^2-1, & C_{\mathrm{Adj}}&=N,\\
  d_{\Lambda^2F}&=\tfrac{N(N-1)}2, & C_{\Lambda^2F}&=\tfrac{(N+1)(N-2)}N,\\
  d_{\mathrm{Sym}^2F}&=\tfrac{N(N+1)}2, & C_{\mathrm{Sym}^2F}&=\tfrac{(N-1)(N+2)}N.
\end{aligned}
\end{equation}

\begin{lemma}[The channel list is complete]\label{lem:channels}
For every $N\ge3$, a second-order process through one shared link has its
intermediate shared-link representation in
$\{\mathbf1,\mathrm{Adj},\Lambda^2F,\mathrm{Sym}^2F\}$ up to conjugation, and
in nothing else. The conjugate of a channel carries the same Casimir and hence
the same weight, which is why four labels suffice where the link-resolved
census of Section~\ref{sec:twocube} needs six.
\end{lemma}

\begin{proof}
The shared link of a one-plaquette shell state carries $F$ or $\bar F$ by
\eqref{eq:EF}; one plaquette action tensors it by $F$ or $\bar F$; and
\eqref{eq:fusion}, together with the conjugate product
$\bar F\otimes\bar F=\Lambda^2\bar F\oplus\mathrm{Sym}^2\bar F$, decomposes
every case completely and multiplicity-free. There is nothing else for the link
to carry.
\end{proof}
\chk{each channel gap is C\_F + C\_R/2, and the weights sum to one}

\noindent The corollary that makes this checkable rather than merely stated is
that the weights $d_\rho/N^2$ sum to one within each orientation family, which
is arithmetic on \eqref{eq:table}: a missing channel would leave a deficit.
What Lemma~\ref{lem:channels} does \emph{not} settle --- and
Section~\ref{sec:scope} says so --- is that one-shared-link processes exhaust
the second-order off-diagonal intermediates. That is a statement about the
process enumeration, not about this list.

The six nonshared half-links contribute $3\CF$ and the fused link $C_\rho/2$,
against an external one-plaquette energy $2\CF$, so the resolvent denominator
is $\CF+C_\rho/2$ exactly.
\chk{each channel gap is C\_F + C\_R/2, and the weights sum to one}

\subsection{The channel weights are a theorem}

Earlier editions assigned the channel projector squared norm $d_\rho/N^2$ by
an isotropy hypothesis, and every statement in this section rests on that
assignment. It is not a hypothesis.

\begin{theorem}[Shared-link channel weights]\label{thm:weingarten}
In the normalised two-plaquette state, the weight of the shared-link channel
$\rho$ is exactly $d_\rho/N^2$, for both orientation families and every
$N\ge3$.
\end{theorem}

\begin{proof}
Two steps. First the six nonshared links collapse. Each plaquette contributes
a product of three independent Haar links, and a product of independent Haar
matrices is Haar, so the amplitude is
$\operatorname{Tr}(AU)\operatorname{Tr}(BU^{\pm1})$ with $A,B$ independent
Haar and $U$ the shared link. Integrating $A$ and $B$ contributes
$\delta/N$ each and leaves a pure degree-$(2,2)$ moment of $U$.

Second, the group. The Weingarten values below are the $U(N)$ ones, and this
is a paper about $SU(N)$; the reduction is a step, not a slip. A degree-$(2,2)$
integrand is invariant under $U\mapsto zU$ for $z\in U(1)$, so integrating it
over $SU(N)$ and over $U(N)=\bigl(SU(N)\times U(1)\bigr)/\mathbb Z_N$ give the
same number provided the $SU(N)$-invariants at this degree are exhausted by the
permutation contractions --- which holds exactly when the degree is below the
rank, $2<N$. That is the theorem's own range, $N\ge3$. At $N=2$ the extra
invariant $\epsilon\otimes\epsilon$ appears and the reduction fails, which is a
second reason the construction begins at $N=3$; it is the same $\epsilon$
channel that makes the $u^1$ row of Appendix~\ref{app:ledger} $SU(3)$-only.
This step has no machine check, and Section~\ref{sec:scope} records it.

Third, two moments settle both families. In the Weingarten calculus,
whose asymptotic origin is~\cite{Weingarten1978} and whose finite-$N$ closed
forms are~\cite{CollinsSniady2006}, the order-two values are
$\mathrm{Wg}(e)=1/(N^2-1)$ and
$\mathrm{Wg}((12))=-1/(N(N^2-1))$ --- the inverse of the $S_2$ Gram matrix
$\bigl(\begin{smallmatrix}N^2&N\\N&N^2\end{smallmatrix}\bigr)$ --- the index
sums give
\begin{equation}\label{eq:moments}
  \begin{aligned}
    M_{\mathrm{direct}}&=\sum_{ijkl}\bigl\langle|U_{ij}|^2|U_{kl}|^2\bigr\rangle=N^2,\\
    M_{\mathrm{cross}}&=\sum_{ijkl}\bigl\langle U_{ij}U_{kl}\bar U_{il}\bar U_{kj}\bigr\rangle=N .
  \end{aligned}
\end{equation}
For the like family, $\operatorname{Tr}(AU)\operatorname{Tr}(BU)$ decomposes
into $\mathrm{Sym}^2\oplus\Lambda^2$ by symmetrising the two index pairs
jointly, and the two squared norms are
$(M_{\mathrm{direct}}\pm M_{\mathrm{cross}})/2$. Normalising,
\[
  n_{\mathrm{Sym}^2}=\frac{N+1}{2N}=\frac{d_{\mathrm{Sym}^2F}}{N^2},
  \qquad
  n_{\Lambda^2}=\frac{N-1}{2N}=\frac{d_{\Lambda^2F}}{N^2}.
\]
For the mixed family the singlet component of $U_{ij}\bar U_{lk}$ is
$\delta_{il}\delta_{jk}/N$, of squared norm $1$ against
$M_{\mathrm{direct}}=N^2$, so $n_{\mathbf 1}=1/N^2=d_{\mathbf 1}/N^2$ and
$n_{\mathrm{Adj}}=1-1/N^2=d_{\mathrm{Adj}}/N^2$.
\end{proof}
\chk{the shared-link weights are Weingarten, not an isotropy assumption}

\begin{remark}
The derivation imports nothing: the Weingarten pair is the inverse of the
$S_2$ Gram matrix and the index sums in \eqref{eq:moments} are explicit, so
the reader can check it without leaving this page. That the pair is not two
fitted constants is separately checked --- the identity-permutation value times
$N^2-1$ is $1$ identically in $N$, so the $SU(3)$ numbers are a specialisation
of a rank-generic formula. Together these are why the all-rank second-order
result is self-contained rather than conditional. Appendix~%
\ref{app:wg} records the index sums.
\chk{the Weingarten route is independent of the corpus}
\end{remark}

\begin{remark}[External corroboration]
The same numerator appears in the quantum-simulation literature from an
unrelated direction. Ciavarella, Burbano and Bauer~\cite{CBB2026} publish a
finite-rank plaquette matrix element
$|M_\rho|^2=\dim\rho/(\dim A\,\dim R)$; for two fundamental factors
$\dim A=\dim R=N$ and this is $d_\rho/N^2$ identically in $N$. That is
provenance for a chain already checked here, not a second derivation, and it
does not check the premise that four channels exhaust the shared-link routes.
\chk{the published dimension-ratio matrix element is this registry's weight formula}
\end{remark}

With Theorem~\ref{thm:weingarten} the resolvent weight is
\begin{equation}\label{eq:w}
  w_\rho=-\frac{d_\rho/N^2}{\CF+C_\rho/2},
\end{equation}
and summing within the two orientation families,
\begin{align}
  A_N&=w_{\mathbf 1}+w_{\mathrm{Adj}}=-\frac{2N^3}{(N^2-1)(2N^2-1)},\label{eq:AN}\\
  B_N&=w_{\Lambda^2}+w_{\mathrm{Sym}^2}=-\frac{4N(N^2-2)}{(N^2-1)(4N^2-9)}.\label{eq:BN}
\end{align}
Each individual weight in closed form is Appendix~\ref{app:closed}.
\chk{the four channel weights follow from dimension and Casimir}
\chk{A\_N and B\_N are the channel sums, not transcriptions}

\subsection{The shared-link law}

The statement that follows is about \emph{one matrix element}. Promoting it to
the whole order-$u^2$ operator on the torus is a separate step with its own
hypotheses, and the two are separated here because the paper's own Scope
section lists one of those hypotheses as unchecked.

\begin{theorem}[All-rank adjacent shared-link coefficient]\label{thm:allrank}
For every integer $N\ge3$, in the untruncated theory or in any regulator
retaining all four fusion channels of \eqref{eq:fusion}, the connected
second-order matrix element between two plaquettes sharing one oriented link is
$t_Nu^2$ times the signed incidence overlap of that pair, with
\[
  t_N=B_N-A_N=\frac{2N(N^2-4)}{(N^2-1)(2N^2-1)(4N^2-9)}>0 .
\]
\end{theorem}

\begin{proof}
Every second-order process between two plaquettes sharing exactly one link runs
through that link's fusion channel; the list is complete by
Lemma~\ref{lem:channels}, the weight of each is Theorem~\ref{thm:weingarten},
and the resolvent denominator is $\CF+C_\rho/2$. Summing within the two
orientation families gives \eqref{eq:AN} and \eqref{eq:BN}, and the relative
orientation of the pair enters as the incidence product, so the two families
appear with opposite sign and the element is $(B_N-A_N)u^2$ times that product.
Positivity: for $N\ge3$ every denominator factor is positive and $N^2-4>0$. The
regulator premise is not decoration --- Corollary~\ref{cor:cutoff} exhibits a
published cutoff under which the same element comes out $-1/12$, of the
opposite sign.
\end{proof}
\chk{t\_N = B\_N - A\_N}\chk{t\_N > 0 for N >= 3}

\begin{conditional}[Incidence assembly on the finite torus]\label{cor:assembly}
Assume in addition that
\emph{(i)} the range of $P_-$ is exactly the charge-odd one-plaquette span, and
\emph{(ii)} after the linked, vacuum-subtracted fold of \eqref{eq:heff} no
connected second-order history between \emph{non-adjacent} plaquettes
contributes a momentum-dependent term on $P_-$.
Then
\[
  H^{(N)}_{\mathrm{eff},-}(k,u)=E_{\mathrm{flat},N}(u)I+t_Nu^2B(k)B(k)\dg+O(u^3),
\]
with $E_{\mathrm{flat},N}$ momentum independent, and the three bands through
second order are $E_{\mathrm{flat},N}(u)$ and $E_{\mathrm{flat},N}(u)+t_Nu^2q(k)$
twice.
\end{conditional}

\begin{proof}
Under \emph{(ii)} only shared-link pairs carry momentum dependence, and
Theorem~\ref{thm:allrank} gives each of those elements, so the off-diagonal
part is $t_NS$ with $S=BB\dg-4I$. The diagonal is a scalar, and for a reason
worth naming because it fails elsewhere in this paper: translations act
transitively on the faces of each orientation and the cubic point group
permutes the three orientations, so the faces are a \emph{single} orbit and any
invariant diagonal is a multiple of the identity. On an open cluster that is
false, and Proposition~\ref{prop:orbit} is what replaces it there. Absorbing
$-4t_N$ into the scalar and applying Theorem~\ref{thm:fact} gives the bands.
\end{proof}
\chk{B B\^{}dagger = q I - d conj(d)\^{}T for the curl incidence}

\noindent Hypothesis \emph{(i)} is the retained-shell premise of
Section~\ref{sec:sector}; \emph{(ii)} is the one
Proposition~\ref{prop:vacuum} narrows but does not close. Neither has a check,
and Section~\ref{sec:scope} says so.

Two consequences,
\begin{equation}\label{eq:tsmall}
  t_3=\frac{5}{612},\qquad
  t_N=\frac1{4N^3}-\frac1{16N^5}-\frac{77}{64N^7}-\frac{1021}{256N^9}+\cdots
\end{equation}
\chk{t\_2 = 0 and t\_3 = 5/612}\chk{large-N expansion of t\_N through 1/N\^{}9}

\subsection{The vacuum-mediated route}

Hypothesis \emph{(ii)} of Corollary~\ref{cor:assembly} is not vacuous, and one
family of routes it has to exclude can be excluded outright. The vacuum-mediated
route connects \emph{any} pair of plaquettes, adjacent or not, and it vanishes
in this sector by a charge-conjugation selection rule. That closes one family
and not the hypothesis; omitting it is a recorded error, not a hypothetical
one.

\begin{proposition}[Vacuum-mediated route]\label{prop:vacuum}
The route $\langle p'|V|0\rangle\langle 0|V|p\rangle/E_0$ is available for any
pair of plaquettes, sharing a link or not. It vanishes identically in the
charge-odd sector, because $V$ is charge-even and $|0\rangle$ is charge-even,
so $\langle0|V|p,-\rangle=0$. Consequently the charge-odd hopping
$t_N=B_N-A_N$ carries no such term, while the charge-even shared-link element
acquires $1/\CF$:
\[
  \ell_N=A_N+B_N+\frac1{\CF}
  =-\frac{2N(3N^2-5)}{(N^2-1)(2N^2-1)(4N^2-9)} .
\]
\end{proposition}
\chk{ell\_N = A\_N + B\_N + 1/C\_F, the vacuum-mediated route at every rank}

\begin{remark}
At $N=3$, $A_3+B_3=-481/612$ and $1/\CF=3/4$, giving $\ell_3=-11/306$. The
first of these is exactly the value once published for the charge-even hopping
and later superseded; the correction is precisely the omitted route. Stating
the selection rule turns a one-rank erratum into an all-rank formula.

One thing this proposition does \emph{not} do is dispose of the route between
non-adjacent pairs. Taken over all pairs it would be a rank-one operator
supported at $k=0$, which is not the twelve-neighbour adjacency $\ell_N$
multiplies; which contributions survive in $H_2$ is settled by the linked,
vacuum-subtracted prescription of \eqref{eq:heff}, not by anything derived
here. What the check above establishes is the all-rank value $\ell_N$ of the
shared-link element, and the disposition of the non-adjacent contributions is
prescription rather than result.
\end{remark}

\subsection{Finite-volume width}

Let $q_{\max}(L)$ be the largest sampled $q$ on the periodic $L^3$ grid. Then
\begin{equation}\label{eq:qmax}
  q_{\max}(L)=
  \begin{cases}
    12, & L\text{ even},\\[2pt]
    12\cos^2\!\bigl(\pi/2L\bigr), & L\text{ odd},
  \end{cases}
\end{equation}
because only at even $L$ does some $k_j$ reach $\pi$ exactly. Under the
hypotheses of Corollary~\ref{cor:assembly}, the width of the
retained charge-odd manifold is $W^{(-)}_{N,L}=q_{\max}(L)\,t_Nu^2+O(u^3)$,
so for even $L$ and in the thermodynamic limit $[u^2]W^{(-)}=12t_N\sim3/N^3$.
This concerns the stiffness of the dispersive doublet, not motion of the
lowest branch, which is flat at this order.
\chk{the zone maximum of q is 12 only at even L}

The smallest nonzero incidence eigenvalue is $q_{\min}=4\sin^2(\pi/L)$,
attained by one unit of momentum in a single direction.
\chk{q\_min on the L-torus grid is 4 sin\^{}2(pi/L)}

\subsection{The orientation signs are essential}

Removing the orientation signs --- from the entries of \eqref{eq:B} and from
the phases, so each $e^{ik}-1$ becomes $e^{ik}+1$ --- gives an unsigned
incidence $\mathcal N(k)$ with $A+4I=\mathcal N\mathcal N\dg$.
The contrast is a determinant identity, not a sample: read straight off the
boundary formulas, $\det B(k)\equiv0$ identically in $k$ --- a zero eigenvalue
at every momentum, which is the whole flat band --- while
$\det\mathcal N(k)=-2\prod_j(1+e^{ik_j})$ vanishes only where some $k_j=\pi$. (Sampling agrees: the unsigned adjacency spectra at $\Gamma$ and $R$
are $\{12,0,0\}$ and $\{-4,-4,-4\}$, which share no eigenvalue.) The flat
branch is therefore not a generic consequence of three face orbitals or cubic
symmetry.
\chk{the two Bloch incidence symbols, and the determinant asymmetry between them}

This control is not hypothetical: the unsigned spectrum is the charge-even
sector's, which Section~\ref{sec:even} computes in closed form, and the two
spans are the two bandwidths. The signed span is $8-(-4)=12$, giving the
charge-odd width $12t_-=5/51$ at $N=3$; the unsigned span is $12-(-4)=16$,
giving the charge-even width $16|\ell_3|=88/153$. Both spans carry the even-$L$
hypothesis of \eqref{eq:qmax}, and for the same reason: the charge-odd $12$
needs every $k_j=\pi$ and the charge-even floor $-4$ needs some $k_j=\pi$, which
the periodic grid samples only when $L$ is even. The charge-even \emph{top}
needs no hypothesis --- it sits at $\Gamma$. Every bandwidth, ratio and band
edge quoted below is therefore an even-$L$, thermodynamic-limit statement unless
it is a $\Gamma$ statement.
\chk{the C-even floor needs even L too, so both bandwidths are even-L statements}
\chk{the two band spans ARE the two incidence spectra}

\section{The second-order coefficient on a two-cube space}
\label{sec:twocube}

Everything above derives $t_N$ from a resolvent ledger indexed by the fusion
channel of one shared link. That ledger is representation theory carried out on
a single plaquette pair; it has never been tested against an operator
calculation on a space large enough for two cubes to interact. This section
reports such a test and, more importantly, states what it settles.

The construction is the open face-sharing $(3,2,2)$ prism with its eleven
oriented plaquettes, in the exact-Casimir $B=6$ $SU(3)$ truncation, of
dimension $1{,}590{,}462$. Its complete charge-odd one-plaquette shell at
$E_*=8/3$ has dimension eleven --- a census result and not the plaquette count.
The full $E_*$ eigenspace has dimension $22$, split $11_{C=+}\oplus11_{C=-}$ by
tensor-derived charge conjugation, and the delivery removes all $22$ from the
pseudoinverse; because the eleven columns exhaust the odd half, that agrees
with the $Q=1_--P_-$ prescription of \eqref{eq:heff}. Writing $K_{LR}$ for the raw second-order
operator and $K_L,K_R,K_F$ for the left-cube, right-cube and shared-face
operators in their own coordinates, the connected operator is defined by the
literal M\"obius fold
\begin{equation}\label{eq:mobius}
  K_{\mathrm{conn}}=K_{LR}-J_LK_LJ_L\dg-J_RK_RJ_R\dg+J_FK_FJ_F\dg ,
\end{equation}
and the same fold applied to the geometry gives $G_{\mathrm{conn}}$. The fold
is literal in the sense that no eigenvalue or fitted scalar is subtracted: each
source operator is folded in its own coordinates and lifted. The delivery
corroborates that internally by applying the same weights history by history
over its $1{,}984$ ordered two-step histories, of which $1{,}192$ cancel at
weight zero, and the two routes agree to $4.9\times10^{-15}$.

First order settles itself. On the shell $PVP=I_{11}$, so the first-order term
is nonzero but scalar: it shifts all eleven branches equally and creates no
branch ambiguity. Its connected fold then vanishes \emph{identically}, and by
counting rather than by cancellation --- each of the eleven faces is covered
exactly once by $\{L,R\}$ once the shared face is restored by $F$, so every
M\"obius weight in \eqref{eq:mobius} applied to the identity is $1-1=0$.
Second order is therefore the leading connected term and not a correction to a
first-order one.
\chk{connected first order vanishes by inclusion-exclusion, not by cancellation of numbers}

At second order the result is
\begin{equation}\label{eq:kconn}
  K^{(2)}_{\mathrm{conn}}=\frac{5}{612}\,G_{\mathrm{conn}}+D ,
\end{equation}
with $D$ diagonal and delivered. The geometry alone fixes the \emph{splitting}
inside each block and nothing else: $G_{\mathrm{conn}}$ turns out to be four
disjoint $2\times2$ blocks plus three isolated directions
(Appendix~\ref{app:twocube}), so each pair splits by $2\times5/612=5/306$ about
its own diagonal entry. Given $D$, which Appendix~\ref{app:twocube} prints,
the spectrum is then closed form,
$\{0,\,(-\tfrac{15}4)^{2},\,(-\tfrac{34}9)^{4},\,(-\tfrac{129}{34})^{4}\}$ ---
exact arithmetic on entries that are themselves reconstructions, a boundary
made precise below.
The delivery reports the coefficient recovered target-blind --- supplied to no
builder, graph fit, branching rule or validator --- and reports that the payload
hashes were produced before the recorded comparison, a chronology it itself
calls record-backed and only partly auditable; the value is byte-present in the
nested $B=4$ package that the $B=6$ build seals as an authority object, so the
defensible statement is dependency-path nonuse and not corpus-wide absence.
This repository and this repository re-thresholds those gates rather
than re-running the build, including a deliberately wrong-sign control that the
held-out scaling rejects: on a $66$-dimensional held-out word space the correct
second-order residual falls with log-log slope $2.90$ while the sign-reversed
control falls with slope $2.00$, a minimum error ratio of $325$. What all this
establishes is runtime nonuse of the comparator on the delivered dependency
path, which is not a preregistration.
\chk{the connected two-cube geometry has exactly four cross-cell pairs, each -1}
\chk{the B=6 connected kernel is (5/612) G\_conn + diag, with the certified spectrum}
\chk{the certificate's own gates pass, target-blind, with the wrong-sign control rejected}

Resolving \eqref{eq:kconn} by the irreducible representation the shared link
actually carries in the ranked intermediate state gives six coefficients rather
than four, because a link distinguishes $\rho$ from $\bar\rho$:
\begin{equation}\label{eq:census}
\begin{array}{c|cccccc}
 \rho & \mathbf1 & \mathbf3 & \bar{\mathbf3} & \mathbf6 & \bar{\mathbf6} & \mathbf8\\
 \hline
 c_\rho & \tfrac1{12} & -\tfrac1{12} & -\tfrac1{12}
        & -\tfrac19 & -\tfrac19 & \tfrac{16}{51}
\end{array}
\end{equation}
\chk{the B=6 six-channel census sums to the registry's own t\_3 = 5/612}

That these six sum to $5/612$ is weak evidence: six rationals reach one target
in many ways, and a fitted set would pass it. The statement worth making is
stronger, and it is what makes this more than a coincidence of one sum. First,
though, the boundary.

The six coefficients of \eqref{eq:census} are not symbolic: the delivery
stores them as rational reconstructions of pinned finite-precision
Clebsch--Gordan contractions, with a maximum reconstruction residual of
$2.2\times10^{-15}$ per channel --- $2.1\times10^{-14}$ across the delivery's
reconstructed matrices as a whole --- and its own certificate declaring
\texttt{symbolic\_exact\_local\_amplitudes: false}. What follows is therefore
reported, conditional on those six numbers, and not a symbolic identity in the
sense of Theorem~\ref{thm:weingarten}.

\begin{reported}[The census is the Weingarten ledger]\label{rep:census}
At $N=3$ the six reconstructed link-irrep coefficients of \eqref{eq:census}
are the four fusion-channel weights \eqref{eq:w} individually:
\[
  c_{\mathbf 1}=-w_{\mathbf 1},\qquad
  c_{\mathbf 8}=-w_{\mathrm{Adj}},
\]
\[
  c_{\mathbf 3}=c_{\bar{\mathbf 3}}=\tfrac12 w_{\Lambda^2F},\qquad
  c_{\mathbf 6}=c_{\bar{\mathbf 6}}=\tfrac12 w_{\mathrm{Sym}^2F}.
\]
\end{reported}

\noindent Substituting \eqref{eq:table} into \eqref{eq:w} at $N=3$ gives
$w_{\mathbf 1}=-1/12$, $w_{\mathrm{Adj}}=-16/51$, $w_{\Lambda^2F}=-1/6$ and
$w_{\mathrm{Sym}^2F}=-2/9$; comparison with \eqref{eq:census} is then
arithmetic. Six reconstructed rationals landing on four predicted ones, with
the two coincidences the correspondence requires, is a strong consistency
statement and not a proof that the underlying contraction is exact.
\chk{the B=6 six-channel census IS the Weingarten four-channel ledger, channel by channel}

One further feature of the census is worth separating out, because it is the
paper's central architecture measured rather than assumed. The six
link-resolved matrices are not merely six numbers that sum correctly: each is
\emph{separately} proportional to the same incidence Gram. The eleven oriented
plaquettes of the prism have exactly $56$ ordered adjacent pairs, each sharing
one link with signed overlap $\pm1$; the delivery reports a constant
channel-to-geometry ratio on all $56$ of them, for every one of the six
channels, with worst residual $2.2\times10^{-15}$. Colour dynamics fixes a
number per channel and the cellular boundary operator fixes the matrix; if the
two mixed, the ratio would be constant for the total and not for the parts.
\chk{every shared-link channel is separately proportional to the geometry, on all 56 pairs}

The correspondence is scoped to what the Casimir closure exhausts: adjacent,
one-action, shared-link intermediates at order $u^2$. Same-face diagonal
routes, on-site terms and higher-action intermediates lie outside it --- the
$B=6$ scalar below is exactly a place where one of them is missing --- and so
does any untruncated-Hilbert-space statement.

It is scoped in one further way, which matters more. The two-cube build draws
its local Wilson amplitudes from the same hash-pinned upstream
Clebsch--Gordan archive as the single-cube reconstructions of the next
subsection, and the squared amplitudes that archive produces are the published
$d_\rho/N^2$~\cite{CBB2026} --- the numerator of \eqref{eq:w}. The census
therefore does not confirm the weights from an independent source: it tests the
\emph{assembly} that carries them, which is the reachable-channel list, the
resolvent denominators, the operator-level fold, the even split between
conjugates, and the geometry that survives each channel separately. That is the
statement worth making, and it is a different statement from confirmation.

Two features of the correspondence carry the content. The mixed family enters
negated because $t_N=B_N-A_N$, which is the same relative sign the incidence
orientation carries. And each like-family \emph{fusion} weight is split evenly
between the irrep and its conjugate --- the two orientations of the shared link
weighted equally --- which the delivery resolves on the two-cube space rather
than assuming. It is not measured \emph{here}: what this repository checks is
the identity the six delivered numbers satisfy.

\begin{conjecture}[All-rank even split]\label{conj:split}
For every $N\ge3$ each like-family fusion weight splits evenly between an
irrep and its conjugate in the link-resolved census.
\end{conjecture}

\noindent This is a conjecture and nothing here promotes it. Falsifying it
means a link-resolved census at some $N\ge4$ in a truncation that still
retains every adjacent shared-link channel; at $N=4$ the endpoint budgets make
that $B\ge33/4$, and no such construction has been run.

\subsection{One truncation, two constructions}

The two-cube space also isolates why a cutoff can reverse the sign of $t_3$.
An adjacent shared-link channel $\rho$ is reachable only if the endpoint
Casimir budget $C_2(\rho)+2C_2(\mathbf 3)$ fits inside the truncation, meaning
$C_2(\rho)+2C_2(\mathbf 3)\le B$: the sextet needs exactly $6$ and the adjoint
$17/3$, so both survive $B=6$ and neither survives $B=4$. The weak inequality is
not a convenience. Both retentions this section uses are decided by an
\emph{equality}: $\Lambda^2F=\bar{\mathbf 3}$ has budget exactly $4$ and so sits
exactly on the $B=4$ cutoff, and the sextet sits exactly on the $B=6$ one. Under
a strict rule $B=4$ would reach $\{\mathbf 1\}$ alone and $B=6$ would lose the
sextet, contradicting both delivered censuses --- so the deliveries are what fix
the convention, which is also why the budget arithmetic is not an independent
confirmation of either census.
\chk{B=6 retains every adjacent shared-link channel; B=4 provably cannot}
\chk{the retention rule is C2(rho) + 2 C2(3) <= B, and both retentions it decides are equalities}

\begin{corollary}[What a link cutoff drops]\label{cor:cutoff}
Retaining only the singlet and $\Lambda^2F$ routes projects the second-order
hopping to
\[
  w_{\Lambda^2F}-w_{\mathbf 1}=-\tfrac1{12},
\]
opposite in sign to $t_3$ and $10.2$ times larger in magnitude. What is
dropped is
\[
  w_{\mathrm{Sym}^2F}-w_{\mathrm{Adj}}=\tfrac{14}{153},
\]
and $-\tfrac1{12}+\tfrac{14}{153}=\tfrac5{612}=t_3$. This is arithmetic on
\eqref{eq:w} alone.
\end{corollary}

\noindent The published $p+q\le1$ link cutoff retains exactly those two
routes --- a per-link restriction, where $B=4$ is a total endpoint budget, so
the two schemes agree on the retained channels at this order and are not the
same rule --- so $-1/12$ is its projected hopping and $14/153$ is the counterterm
that completes it --- not $5/612$, which would double-count what it already
carries. On the two-cube space the first of those two numbers is delivered
rather than inferred: a separately sealed $B=4$ run on the identical geometry
gives $K^{(2),B4}_{\mathrm{conn}}=-\tfrac1{12}G_{\mathrm{conn}}+D_{B4}$
outright, and the Casimir budgets of the previous paragraph say without
reference to any census that $B=4$ can reach only
$\{\mathbf1,\mathbf3,\bar{\mathbf3}\}$. Conditional on \eqref{eq:census}, the
same split of the census reproduces it,
$c_{\mathbf 1}+c_{\mathbf 3}+c_{\bar{\mathbf 3}}=-1/12$, and gives the omitted
half as $c_{\mathbf 6}+c_{\bar{\mathbf 6}}+c_{\mathbf 8}=14/153$.
\chk{the B=6 six-channel census IS the Weingarten four-channel ledger, channel by channel}
\chk{FINDING: the T1 link cutoff reverses the sign of t\_3, and 14/153 is what it omits}

\begin{remark}
Two single-cube reconstructions bracket the same statement. A full $B=4$ open
cube assembled from the published \texttt{ymcirc} $d=3$, $B=4$ plaquette
table~\cite{Balaji2026} reproduces $-1/12$ and the truncated shell ordering to
$4.3\times10^{-11}$; a $B=6$ cube reproduces $+5/612$ and the inverse ordering
to $8.4\times10^{-15}$. Both locate the missing weight at the same $14/153$.

Neither is an independent replication. Both read tabulated plaquette matrix
elements of the same lineage --- and the two-cube build draws its
Clebsch--Gordan amplitudes from the same hash-pinned \texttt{pyclebsch} archive
as the $B=6$ cube, so that agreement is a cross-check of assembly and not of
the amplitudes. The published closed form~\cite{CBB2026} is
$d_\rho/N^2$ --- the very numerator of \eqref{eq:w}. This is a cross-check of
assembly, not a third confirmation of $t_3$; the exact rational identity is
Corollary~\ref{cor:cutoff}, which needs no tolerance. The $p+q\le1$
nomenclature and the finite-rank matrix element are both from~\cite{CBB2026},
and single-cube $SU(N)$ diagonalisation in this style goes back
to~\cite{CM2003}.

A counterterm added to a published $T_1$ Hamiltonian is therefore $14/153$ and
not $5/612$, which would double-count the two routes $T_1$ already carries. One coincidence of value is worth naming:
$-1/12$ is also the singlet weight $w_{\mathbf 1}$, and the cutoff value is a
\emph{difference} of two weights, not a weight.
\chk{FINDING: a full T1 = B = 4 cube Hamiltonian reproduces -1/12 and the reversed shell}
\chk{FINDING: the B = 6 cube flips the sign back to +5/612, closing the decisive test}
\end{remark}

The $B=6$ cube is channel-complete for the hopping and \emph{not} for the
scalar: its scalar is $39/68$ against the bridge's $11/34$, differing by
exactly $1/4$, the local same-face sextet route, which first enters at cutoff
$20/3$ --- integer label $B=7$, so any cutoff above $20/3$ already carries it.
That is not only arithmetic: the delivery reports the prediction executed. A
six-face $B=7$ reachable-state probe returns $\alpha_{B=7}=11/34$ and
$t_{B=7}=5/612$ with scalar-plus-Gram residual $1.1\times10^{-15}$, so raising
the cutoff past the same-face threshold lowers all three absolute coefficients
by exactly $1/4$ and leaves the relative shell $\{0,\,5/153,\,5/102\}$
untouched --- the scalar moves, the hopping does not. That probe's artifact did
not travel, so it is reported here and not checked.
\chk{the B = 6 scalar misses the bridge's by exactly the same-face sextet route}

\subsection{What the cutoff moves, and where}

The connected diagonal $D$ is eleven numbers in both deliveries and is treated
as a remainder in both. It is not eleven numbers. The open prism has a symmetry
group --- the cell exchange $x\mapsto2-x$, the two transverse reflections and
the $y\leftrightarrow z$ swap, of order $16$ --- and $D$ is constant on its
orbits, which are $8+2+1$: the eight side faces, the two end caps, the shared
face.

\begin{proposition}[Where the truncation lives]\label{prop:orbit}
The eight-element orbit is exactly the support of the four cross-cell pairs,
and it is the only orbit whose diagonal value differs between the two
truncations: the end caps carry $-15/4$ and the shared face $0$ at both $B=4$
and $B=6$, while the side value moves from $-7/4$ to $-2317/612$. Together with
the four cross-cell blocks carrying the whole off-diagonal change, the entire
$B=4\to B=6$ difference is confined to the faces that carry connected transport.
\end{proposition}
\chk{the connected diagonal is orbit-constant, and only the transporting orbit moves}

This also says what the nonuniform $D$ is. Where the faces of a complex form a
\emph{single} orbit --- the closed cube of the next subsection, the periodic
torus of Theorem~\ref{thm:allrank} --- orbit-constancy forces the diagonal to
be scalar, and the operator is exactly $\alpha I+tG$. The prism has three
orbits because it has a boundary. $D$ is the open boundary showing up in the
diagonal, not a failure of the incidence form.

\subsection{The one-cube shell, and why its flat level cannot move}

Those two single-cube reconstructions are usually read as numerical
cross-checks. They are more than that, because the object they diagonalise is
the same chain complex. The six faces of a cube are a closed surface, so in the
outward orientation $\ker\partial_2$ is one-dimensional and spanned by the sum
of all six --- the fundamental class, $b_2(S^2)=1$. That is
Theorem~\ref{thm:carrier} on a complex with no three-cells, where the $L^3-1$
boundaries are absent and only the class survives.

\begin{proposition}[One-cube shell]\label{prop:cube}
On the closed cube surface $G=\partial_2^{\mathsf T}\partial_2$ commutes with
all $24$ proper rotations. The combinations $b_{+i}-b_{-i}$ carry the defining
vector representation and the combinations $b_{+i}+b_{-i}$ the permutation
representation of the three unordered axes, so with the charge-odd inversion
$Pb_{+i}=-b_{-i}$ the shell is
\[
  A_1^{--}\oplus T_1^{+-}\oplus E^{--},
\]
with $G$-eigenvalues $0$, $4$, $6$ and multiplicities $1$, $3$, $2$.
Consequently $\alpha I+tG$ has levels $\alpha$, $\alpha+4t$, $\alpha+6t$.
\end{proposition}

\begin{proof}
The rotations permute the six outward faces, and $G$ is built from the boundary
map, so it commutes with the permutation action. On the antisymmetric
combinations that action is the defining representation by inspection --- $b_{+i}-b_{-i}$
transforms as $e_i$ --- and on the symmetric combinations it is the
permutation representation of $\{\pm e_i\}$, which is $A_1\oplus E$. Under
$Pb_{+i}=-b_{-i}$ the antisymmetric sector is parity even and the symmetric one
parity odd. In these sectors $G$ is $4I$ and $4I-2(J-I)$ respectively, of
spectra $\{4,4,4\}$ and $\{0,6,6\}$; since the three isotypic components have
distinct dimensions the eigenvalues attach to them uniquely.
\end{proof}
\chk{the one-cube shell is A1 + T1 + E, and its flat level is the S\^{}2 fundamental class}

Three things follow. The cubic labels are derived here rather than assigned:
both deliveries assert them and neither builds the symmetry matrices. The
$A_1^{--}$ level is Proposition~\ref{prop:rigid}'s mechanism on a second
complex --- $Gw=0$ there, so that level sits at $\alpha$ for every $t$, and no truncation, no
channel and no coefficient in dispute anywhere in this paper can move it. And
the reversal the two runs report is therefore a statement about the sign of one
number rather than about the shell as a whole: at $B=6$ the levels are
$\{39/68,\,371/612,\,127/204\}$ and at $B=4$ $\{37/12,\,11/4,\,31/12\}$, in
the opposite order, with $A_1^{--}$ in each case sitting exactly at the scalar.
The relative shell $\{0,4t,6t\}$ is $\{0,\,5/153,\,5/102\}$ at $B=6$.

The finite-volume fingerprint of the 28 August finite-$N$ bridge certificate
--- a different release, and named here because a reader tracing it to the
two-cube one would find nothing --- is worth recording
because its most discriminating multiplicity is not a feature of the
truncation: on the periodic $3^3$ torus the incidence spectrum is
$\{0^{29},3^{12},6^{24},9^{16}\}$, and $29=L^3+2$ is
Theorem~\ref{thm:carrier}.
\chk{the certificate's finite-volume fingerprints, and 29 = L\^{}3 + 2 is the Lean cycle count}

\begin{remark}[Scope]
Everything in this section is second order in $u$, finite volume, strong
coupling, one sector. It bears on $t_N$ and decides nothing about the
fourth-order coefficient of Section~\ref{sec:fourth}. The
$B=6$ two-cube construction was not re-run here: its arithmetic is
audited against geometry rebuilt independently, and its diagonal $D$ is
$B$-truncated and not proved cutoff-stable. The held-out finite-$u$ validation
that rejects the wrong-sign control is a $66$-dimensional $P+Q_1$ word space
with the $Q_1$--$Q_1$ magnetic block set to zero, not the full Hamiltonian, and
no radius-three stability claim was executed. The two branches are not equally
covered: at $B=4$ the delivery reports a direct diagonalisation of the complete
$4{,}180$-dimensional charge-odd block at four small couplings, two of them
held out, with the residual after $E_*+u+u^2\operatorname{eig}K_2$ scaling as
$u^3$; at $B=6$ no such diagonalisation exists and the $66$-dimensional star
stands in for it. That $B=4$ report carries no check here --- its certificate
did not travel --- and is recorded as the delivery's rather than as this
paper's. No external group has reproduced the release.
\end{remark}

\section{The charge-even sector in closed form}
\label{sec:even}

The unsigned incidence $\mathcal N(k)$ of Section~\ref{sec:second} is not a
hypothetical control: it is the charge-even sector's Bloch symbol. Write
\[
  \hat a_m=2+2\cos k_m=|1+e^{ik_m}|^2,\qquad p=\sum_m \hat a_m .
\]
The hat is not decoration: $\hat a_m$ is the \emph{complement} of the
$a_m=4\sin^2(k_m/2)$ used in the charge-odd shape basis \eqref{eq:shapes}
below,
$\hat a_m+a_m=4$, and that complementarity is the sum rule below.

\begin{theorem}[Charge-even Bloch cubic]\label{thm:evencubic}
$p+q=12$ pointwise, and the characteristic polynomial of
$\mathcal N\mathcal N\dg$ is
\[
  \mu^3-2p\,\mu^2+p^2\mu-4\hat a_1\hat a_2\hat a_3
  \;=\;\mu(\mu-p)^2-4\hat a_1\hat a_2\hat a_3 ,
\]
with the charge-even adjacency eigenvalues $\lambda=\mu-4$. The charge-odd
counterpart of the same computation is $\mu(\mu-q)^2$.
\end{theorem}
\chk{the C-even characteristic polynomial is mu(mu - p)\^{}2 = 4 a\_1 a\_2 a\_3}
\chk{p + q = 12: one zone function runs both sectors}

The closed form is complete rather than merely consistent at sampled momenta:
the characteristic polynomial of the $81\times81$ and $192\times192$ finite
plaquette adjacencies factorises coefficient for coefficient over $\mathbb Z$
into this cubic over the momentum grid.
\chk{the Bloch cubic IS the finite L = 3 and L = 4 plaquette spectrum, exactly}

\begin{proposition}[Range and attainment]\label{prop:evenrange}
$\lambda\in[-4,12]$ exactly. The top $\lambda=12$ is attained only at $\Gamma$
($\hat a_1=\hat a_2=\hat a_3=4$). The bottom $\lambda=-4$ is attained on the
whole of the three planes $k_j=\pi$, where the constant term vanishes, and only
at $R$ is it a triple root.
\end{proposition}

\begin{proof}
Write $f(\mu)=\mu(\mu-p)^2-4\hat a_1\hat a_2\hat a_3$. Then
$f(0)=-4\hat a_1\hat a_2\hat a_3\le0$ and $\mu(\mu-p)^2<0$ for $\mu<0$, so no
root lies below $0$; and $f'(\mu)=(3\mu-p)(\mu-p)>0$ for $\mu>p$, with
$f(16)=16(16-p)^2-4\hat a_1\hat a_2\hat a_3\ge0$ because $p\le12$ and
$\hat a_1\hat a_2\hat a_3\le64$. Hence $\mu\in[0,16]$ and
$\lambda=\mu-4\in[-4,12]$. Equality at the top forces $p=12$ and
$\hat a_1\hat a_2\hat a_3=64$, which is $\Gamma$. Equality at the bottom is
$f(0)=0$, i.e. $\hat a_1\hat a_2\hat a_3=0$, i.e. some $k_j=\pi$; there
$f(\mu)=\mu(\mu-p)^2$ and the root $\mu=0$ is simple unless $p=0$, which is
$R$ alone.
\end{proof}
\chk{the C-even range [-4, 12] is exact; the top only at Gamma, the floor on three planes}
\chk{the C-even band touches its floor exactly on the three planes k\_j = pi}

\begin{remark}
A flat charge-even branch would need $f(\mu_0)=0$ identically in $k$, which the
constant term already forbids. The charge-even kernel is three measure-zero
planes against the charge-odd kernel's whole band: the two sectors are two
readings of one geometric object, and only the signed one is flat.
\chk{the C-even band touches its floor exactly on the three planes k\_j = pi}
\chk{the C-even spectra at the four high-symmetry momenta}
\end{remark}

\subsection{One assembly, both sectors, both orders}

The $r=2$ content below is all-rank and checked. The $r=3$ values ---
$t_{3,+}=-6335/249696$ and everything computed from it --- come from the same
supplied domino ledger as Section~\ref{sec:su3}, and are conditional on it in
the same way.

The scalar ledger is not a list of separate facts. Writing $\lambda$ for the
adjacency eigenvalue and $\mathrm{tower}_{r,s}$ for the within-plaquette term
in canonical $u$,
\begin{equation}\label{eq:assembly}
  E_s(\lambda,r)=\mathrm{tower}_{r,s}+12\,\mathrm{leak}_{r,s}+\lambda\,t_{r,s},
\end{equation}
with $\lambda\in\{-4,8\}$ in the charge-odd sector (carrier and dispersive
top), $\lambda\in\{12,-4\}$ in the charge-even sector, and $\lambda=0$
returning the on-site diagonals. This reproduces every registered diagonal
and band value --- nine of them --- across both sectors and both orders. The
tower term is the certified coupling conversion $4\Delta(3u/2)$, so
\eqref{eq:assembly} takes no unregistered input. The twelve is a count on this
lattice, not a citation: the plaquette graph is $12$-regular and two distinct
faces share at most one link. The two towers are different objects and both
are needed: the charge-odd one contributes
$8/3+u+u^2/2+7u^3/32$ and the charge-even one $8/3-u+(13/20)u^2+(101/200)u^3$,
each of them $4\Delta(3u/2)$ on its own sector's series, so
$\mathrm{tower}_{2,-}=1/2$ against $\mathrm{tower}_{2,+}=13/20$ and
$\mathrm{tower}_{3,-}=7/32$ against $\mathrm{tower}_{3,+}=101/200$.
The nine values it reproduces are tabulated in Appendices~\ref{app:ledger}
and~\ref{app:even}.
\chk{one assembly formula gives every registered band value}
\chk{one assembly formula gives every C-even value at both orders}
\chk{the plaquette graph is 12-regular and two faces share at most one link}

Because $\mathrm{leak}_{r,+}=t_{r,+}$ at both computed orders, the whole
charge-even sector collapses to
$\mathrm{tower}_{r,+}+(12+\lambda)\,t_{r,+}$, a one-parameter family in the
charge-even tower. That
equality is verified and unexplained. The order-three separation is between the
charge-odd leakage and the charge-\emph{even} hopping --- not, as an earlier
edition printed, between the charge-odd leakage and the charge-odd hopping,
which were never equal: $\mathrm{leak}_{2,-}=t_{2,+}=-11/306$ against
$t_{2,-}=5/612$ at second order, and $\mathrm{leak}_{3,-}=-12331/249696$ against
$t_{3,+}=-6335/249696$ at third. It is recorded as a unifying candidate with a
falsifier --- an order at which the two are computed independently and differ
--- not as a mechanism.
\chk{both declared coincidences, checked: one is ell\_N at all ranks, the other is bare}

\subsection{Where the rest frame is not blind}

\begin{proposition}[Charge-even $\Gamma$ splitting]\label{prop:evengamma}
At $\Gamma$ the charge-even spectrum is $\{12,0,0\}$, so the $A_1^{++}$ level
and the $E^{++}$ doublet are distinct, and
\[
  E(A_1^{++})-E(E^{++})=12\,t_{r,+}
\]
exactly: $-22/51$ at $r=2$ and $-6335/20808$ at $r=3$.
\end{proposition}
\chk{the C-even Gamma point pins t\_+, exactly where no C-odd Gamma datum can}

This is the structural counterpart of the charge-odd blindness recorded in
Section~\ref{sec:su3}. The signed spectrum at $\Gamma$ is $\{-4,-4,-4\}$, one
level, so no charge-odd rest-frame series can see the hopping; the unsigned
spectrum has two levels, so the charge-even rest frame can. Neither half needs
the incidence ansatz, and stating them through it understates both. At $\Gamma$
the little group is the full cubic point group, so the $3\times3$ block commutes
with its action at every order in $u$ and whether or not the correction factors
through links; Schur then decides the two sectors, and decides them
differently. The charge-odd state is the \emph{oriented} $2$-cell
$e_i\wedge e_j$, on which a proper rotation acts by its exterior square, which
is the rotation itself under Hodge duality --- the defining rep $T_1$,
irreducible, commutant dimension one, so the block is a scalar. Charge
conjugation drops the orientation, so the charge-even state is the
\emph{unoriented} plane and the action is the permutation of the three axes,
$A_1\oplus E$, commutant dimension two. The two commutant dimensions are
exactly the two level counts. The $A_1^{++}$
branch is isotropic at leading order without assuming it, with curvature
$\tfrac43|t_{r,+}|$ in all three directions --- $22/459$ at second order and
$6335/187272$ at third.
\chk{the C-even curvature is (4/3)|t\_+| at both orders, and isotropic}
\chk{Schur alone fixes both Gamma blocks: the C-odd triple is irreducible, the C-even one is not}

At both computed orders the charge-even bandwidth is $16|t_{r,+}|$ and the
charge-odd manifold width $12|t_{r,-}|$. At any order at which the correction
stays in the incidence algebra --- which Proposition~\ref{prop:rigid} warns is
not automatic --- the second-order ratio is $4(3N^2-5)/[3(N^2-4)]$, finite and
greater than one at every rank $N\ge3$, decreasing from $88/15$ at $N=3$
towards $4$ --- a ratio of two even-$L$ spans, in the sense fixed above. So at second order the charge-even band is the wider one at every
rank, and that asymmetry is a channel statement, not a geometric one.
\chk{the C-even bandwidth is 16|t\_+| at every order; the C-odd manifold width is 12|t\_-|}
\chk{at N = 2 the C-odd hopping vanishes and the C-even one does not}

\section{$SU(3)$ through third order}
\label{sec:su3}

The \emph{third-order} content of this section is conditional on supplied
exact-arithmetic ledgers. Its second-order content is not: $\mathrm{leak}_2$ is
the all-rank $\ell_N$ of Proposition~\ref{prop:vacuum} at $N=3$, and the tower
is the certified coupling conversion. The within-plaquette charge-odd tower
contributes $8/3+u+u^2/2+7u^3/32$, and the registered $SU(3)$ values, with the
scope of each row, are collected in Appendix~\ref{app:ledger}.

\begin{reported}[Third-order ledger]\label{rep:su3}
The supplied ledgers report, in exact rationals,
\[
  b_3=\frac{1975}{124848},\qquad
  \mathrm{leak}_3=-\frac{12331}{249696},
\]
and the assembled flat scalar
$d_3=\tfrac7{32}+12\,\mathrm{leak}_3-4b_3=-\tfrac{109151}{249696}$.
\end{reported}
\chk{d\_3 = 7/32 + 12*leak\_3 - 4*b\_3}
\chk{leak\_3 is assembled from the domino diagonal and the vacuum piece}

\begin{conditional}[Third-order factorisation]\label{cor:su3}
If the lifter census is exhaustive and the linked ledger complete, then
\[
  H_{\mathrm{eff},-}(k,u)=E_{\mathrm{flat}}(u)I+\tau(u)B(k)B(k)\dg+O(u^4),
\]
\[
  \begin{aligned}
    E_{\mathrm{flat}}(u)&=\tfrac83+u+\tfrac{11}{306}u^2-\tfrac{109151}{249696}u^3,\\
    \tau(u)&=\tfrac5{612}u^2+\tfrac{1975}{124848}u^3 .
  \end{aligned}
\]
\end{conditional}
\chk{E\_flat and t(u) carry the ledger coefficients}

The second-order diagonal is the $\lambda=0$ value
$\tfrac12+12\,\mathrm{leak}_2=7/102$ with $\mathrm{leak}_2=-11/306$;
subtracting $4t_-$ gives the flat scalar $11/306$ at $\lambda=-4$.
\chk{d\_- = 1/2 + 12*leak\_2, and leak\_2 = -11/306}

\begin{remark}[What the external comparison can and cannot do]
Hamer's 1989 rest-frame series~\cite{Hamer1989} agrees with $E_{\mathrm{flat}}$
at every shared order. It cannot test $\tau$. The hopping enters the spectrum
only through $\tau(u)q(k)$ and $q(0)=0$, so every \emph{charge-odd}
$\Gamma$-point datum is blind to it. That argument runs inside the incidence
ansatz, which Proposition~\ref{prop:rigid} warns is not unconditional; the
Schur argument of Section~\ref{sec:even} removes the caveat, since the
charge-odd triple carries the irreducible $T_1$ and its $\Gamma$ block is
therefore scalar at every order. The charge-even rest frame is not blind, by
Proposition~\ref{prop:evengamma}; the one-parameter family
$(b_3+\varepsilon,\ \mathrm{leak}_3+\varepsilon/3)$ leaves $d_3$ and therefore
the entire comparison unchanged. What the agreement pins is the combination
$12\,\mathrm{leak}_3-4b_3$. The $\lambda=8$ band top separates them, via \eqref{eq:assembly}.
Proposition~\ref{prop:evengamma} is the charge-\emph{even} analogue and
separates $\mathrm{leak}_+$ from $t_+$; it says nothing about $b_3$ or
$\mathrm{leak}_3$.
\chk{FINDING: no Gamma-point datum can constrain the hopping}
\end{remark}

\section{What the homology protects}

\begin{proposition}[Boundary-factorised rigidity]\label{prop:rigid}
If an order-$r$ Bloch correction has the form
$H_r(k)=a_rI+b_rS(k)+B(k)M_r(k)B(k)\dg$, then for every $k\neq0$,
$H_r(k)w(k)=(a_r-4b_r)w(k)$: the correction shifts the carrier without
dispersing it.
\end{proposition}

\begin{proof}
$B\dg w=0$ and $Sw=-4w$, both from Theorem~\ref{thm:fact}. The middle term is
annihilated whatever $M_r$ is, which is the content: the protection belongs to
the boundary operator and not to the dynamics.
\end{proof}
\chk{a boundary-factorised correction shifts the carrier without dispersing it, for generic M}

In real space this is the orthogonality of the two cellular Laplacians,
$L^{\downarrow}_2L^{\uparrow}_2=L^{\uparrow}_2L^{\downarrow}_2=0$, so any
polynomial in them acts on $\ker L^{\downarrow}_2\cap\ker L^{\uparrow}_2$ by
its constant term. The proposition is not an unconditional all-orders
statement: $\{BMB\dg\}$ is not a two-sided ideal in $\mathrm{End}(C_2)$, and a
general higher-order correction need not factor through links. The
carrier-projected fourth-order shapes the corpus works with are
\begin{equation}\label{eq:shapes}
  1,\quad q,\quad e_2,\quad \frac{4e_2}{q},\quad \frac{e_3}{q},
\end{equation}
with $e_2=\sum_{i<j}a_ia_j$, $e_3=a_1a_2a_3$ and $a_i=4\sin^2(k_i/2)$. The
$1/q$ is the carrier projection: $\|w\|^2=q$ by Theorem~\ref{thm:fact}, so the
carrier-projected correction is $\varepsilon_4=(w\dg H_4w)/q$ and the
numerator, not $\varepsilon_4$, is the polynomial.
\chk{the carrier projection is where the 1/q comes from}
\chk{clearing the denominator reproduces the five-element numerator basis}

That family is \emph{not} what cubic symmetry permits, and an earlier edition
of this paper attributed it there. The point group permutes $(a_1,a_2,a_3)$, so
its invariants are the symmetric polynomials, and the numerators of degree at
most three with no constant term span a six-dimensional space
$\{q,q^2,e_2,q^3,qe_2,e_3\}$ --- one per partition. Clearing the denominator in
\eqref{eq:shapes} returns five of those six; the missing direction is $q^3$,
that is the shape $q^2$, which is cubic invariant and regular at $\Gamma$
($q^2\sim|k|^4$), so neither the point group nor the carrier projection
excludes it. The omission is a property of the two-hop enumeration that
produced the family, and it is load-bearing: Theorem~\ref{thm:obstruction} is a
rank statement about \emph{this} span, and a sixth direction would be a sixth
unknown for the same data to determine.
\chk{the shape family is not symmetry-complete: cubic symmetry alone permits q\^{}2 as well}

\section{The fourth-order boundary}
\label{sec:fourth}

Two supplied fourth-order records agree on the axial datum and on the
appearance of mobility, and disagree on one off-axis mixed-gradient
coefficient $C_{\mathrm{shp}}$. This edition does not adjudicate them, and no
fourth-order band or bandwidth is used above. What it does is state precisely
where the disagreement lives, because that geography is now checked and it is
what a decisive calculation must target.

The two values, side by side and in no order of preference:
\begin{equation}\label{eq:c2}
\begin{aligned}
  C_{\mathrm{old}}&=-\frac{211835444920651}{4405310420659200}=-0.0480863\ldots,\\
  C_{\mathrm{new}}&=-0.020213328886166577,\\
  \Delta_C&=C_{\mathrm{new}}-C_{\mathrm{old}}=0.0278730\ldots
\end{aligned}
\end{equation}
The first is recorded as an exact rational and the second only as a float ---
that asymmetry is a fact about the two records, not an argument for either. The
gap is what ``less than half the disputed gap'' and ``the two records give
$K=17.04$ and $K=23.66$'' below are measured against.
\chk{Delta\_C = C\_new - C\_old > 0 (the real discrepancy)}

The two records do agree on the axial datum, but the corpus's own quoted
tolerance does not cover that agreement: the $\alpha$ gap is
$2.43\times10^{-13}$ against a stated bound of $2.3\times10^{-13}$. The sealed core holds internally --- $\alpha_{\mathrm{new}}$ and
$4A_{\mathrm{new}}$ agree to $2\times10^{-15}$ --- which is consistent with the
quoted bound being too tight, but two float comparisons do not discriminate
that from a shared upstream slip. Recorded as a discrepancy rather than
resolved, and the bound is not widened.
\chk{v10a.26 A, B, D match the sealed rationals within 2.3e-13}
\chk{FINDING: alpha\_new falls outside the corpus's own 2.3e-13 bound}

\subsection{The obstruction is polynomial}

\begin{theorem}[Obstruction certificate]\label{thm:obstruction}
The two recorded off-axis coefficients enter \eqref{eq:shapes} only through
the shape $4e_2/q$, so the two records' carrier numerators differ by
$4\Delta_C e_2$ and their bands by $\Delta_C\,(4e_2/q)$. On any
one-dimensional axial cut two of the three $a_i$ vanish, so $e_2$ is the
\emph{zero polynomial} there and both differences vanish identically.
Consequently no $\Gamma$-point or axial datum distinguishes the two records,
at any precision. Off axis they do differ: $e_2(M)=16$ and $e_2(R)=48$, so the
numerators differ by $64\Delta_C$ and $192\Delta_C$ and the bands by
$8\Delta_C$ and $16\Delta_C$.
\end{theorem}
\chk{FINDING: the retained Gamma/axis data cannot identify C\_shp}
\chk{Phi\_C vanishes at Gamma along every direction}

\begin{remark}
The vanishing is polynomial, not numerical: this is non-identifiability, not a
limitation of the data in hand. Re-anchoring cannot substitute, and neither
record is preferred here. An off-axis contraction is the only decider.
\chk{the crosswalk is exactly scalar on the momentum axes}
\end{remark}

On the axial cut itself the invariants collapse to $q=L(k)=4\sin^2(k/2)$ with
$e_2=e_3=0$, and dividing the raw cube-boundary numerator $(\alpha_N/4)L^2$ by
$\|w\|^2=q=L$ leaves a single power of $L$ with coefficient $\alpha_3/4=5/48$.
\chk{on an axial cut the mixed invariants vanish and the norm divides}

\subsection{A second witness}

Non-identifiability is better exhibited than argued from a difference between
two records, since a difference could always be a computational error in one of
them. The exhibition adds no logical strength: Theorem~\ref{thm:obstruction}
already proves that \emph{every} value of $C$ is consistent with the retained
data, so a witness is one point of a line already known to be a line. It is a
concrete handle, not evidence. The witness below is not itself a candidate for
the physical coefficient --- it is the balanced $\epsilon$-free continuation of the
all-rank shape family, and $\epsilon$-free is precisely what the physical value
is not. Its job is to make the $4\Delta_C e_2$ family concrete with an exactly
known shift.

\begin{proposition}[Explicit second witness]\label{prop:witness}
There is a value $C_{\mathrm{alt}}$ with
$C_{\mathrm{alt}}-C_{\mathrm{old}}=25/1024$ exactly, appearing verbatim in the
pinned source edition, whose induced band separations are $0$ at $X$,
$25/128$ at $M$ and $25/64$ at $R$. It is identical to
$C_{\mathrm{old}}$ on every retained datum and distinct from it off-axis.
\end{proposition}
\chk{FINDING: an explicit second witness C\_alt exhibits the C2 non-identifiability}
\chk{the C\_alt witness IS the balanced eps-free continuation, and Delta\_C/(A/2) = 15/32}

\subsection{Where the coefficient lives}

The four checkpoint deltas are the $\Gamma$-subtracted carrier corrections
$\Delta_s=\varepsilon_4(k_s)-\varepsilon_4(0)$ at
$k_X=(\pi,0,0)$, $k_M=(\pi,\pi,0)$, $k_P=(\pi,\pi/2,0)$ and
$k_R=(\pi,\pi,\pi)$ --- $P$ is the off-axis, off-corner point that makes the
system invertible, with $a=(4,2,0)$. In the shape basis they are
\begin{equation}\label{eq:checkpoints}
\begin{aligned}
  \Delta_X&=4A,\\
  \Delta_M&=8A+16B+8C,\\
  \Delta_P&=6A+8B+\tfrac{16}3C,\\
  \Delta_R&=12A+48B+16C+\tfrac{16}3D,
\end{aligned}
\end{equation}
and they invert exactly, so the four shape coefficients are recoverable from
four checkpoints. $\Delta_X$ contains $A$ alone, which is why the axial cuts
agree whatever the dispute does.
\chk{checkpoint values at X, M, P, R}
\chk{the four extraction formulas invert the ansatz}
\chk{X is blind to B, C, D — it fixes A alone}

Four further measurements locate the coefficient and bound where a
disagreement could live. The distinction matters and is easy to lose: showing
where a coefficient's \emph{value} is concentrated is not by itself showing
where a \emph{discrepancy} can sit. Only the third item below does the second
thing.

\begin{enumerate}\itemsep2pt
\item \textbf{The shape fit's $C$ row sums to zero}, so no error in the
  $\Gamma$ anchoring can move $C_{\mathrm{shp}}$ at all; the recorded run's own
  re-anchoring moved it by $4.6\times10^{-15}$, as the zero row sum requires.
  A scalar shift is translation-local and changes nothing observable, which is
  why the two anchored scalars $q_{\mathrm{band}}^{(4)}$ and
  $m_\Gamma^{(4)}$ are differently anchored coordinates and not rival
  estimates.
\item \textbf{The coefficient's value is concentrated, not spread.} Six of
  the $189$ kernel records --- the normal $(0,0,1)$ shell --- carry $81\%$ of
  the total weight, while the $120$ rotation records carry $6\%$ and are $255$
  times less sensitive per record. This says where the coefficient lives, not
  yet where the two records can differ.
\item \textbf{The agreed axial coefficient pins the normal sector.} All six
  normal records share one amplitude $h$, and it sets both shapes rigidly:
  $A_{\mathrm{normal}}=-h$ while the raw normal amplitude is $2h$, which in the
  shape normalisation of \eqref{eq:shapes} is $C_{\mathrm{normal}}=h/2$, so
  \[
    C_{\mathrm{normal}}=-\tfrac12 A_{\mathrm{normal}}\quad\text{exactly.}
  \]
  (The factor of four between $2h$ and $h/2$ is the raw-to-shape conversion,
  not a discrepancy.) Granting $A=5/48$ --- which both records share --- together
  with the block decomposition, which puts $5/48+4x$ on the normal block and
  $-4x$ on the mixed one, therefore fixes $h$ and with it $81\%$ of the
  coefficient; and everything free to move without disturbing $A$ totals less
  than half the disputed gap. A rival kernel must differ \emph{structurally} in
  the $A$-carrying sector; no redistribution reaches the gap.
\item \textbf{The two kernels agree almost everywhere.} Their supports are
  identical and $144$ of $189$ records ride one common scale; the divergent
  cross sector is a single uniform ratio. The dispute is three amplitudes.
\end{enumerate}
\chk{the shape fit's C row sums to zero, so no Gamma-anchor error can move C\_shp at all}
\chk{a translation-local scalar shift changes nothing observable}
\chk{FINDING: C\_shp is carried by 6 of 189 records, not spread across the kernel}
\chk{FINDING: A = 5/48 pins the normal sector's whole C4 contribution}
\chk{C\_normal = -A\_normal/2: the agreed axial coefficient pins the normal channel}
\chk{the two 189-record kernels agree everywhere except three amplitudes, and the on-site anchor swap moves C by exactly zero}

Read together these four sharpen the target and prefer no side of it. They say
that a rival kernel must differ structurally in the $A$-carrying sector rather
than by redistribution. That is a constraint on candidates built on the
recorded $189$-record support --- which is where the measurement was made; a
kernel on a different support is not constrained by it at all. Neither record is promoted here and neither is withdrawn: the interval of
\emph{values} between them is not narrowed. What is narrowed is the class of
\emph{kernels} that can reach either end, and that constraint bears on both
records alike.

\subsection{What cannot decide it}

A large sealed sweep over $609$ marked clusters has been the standing candidate
for settling this. It cannot: its engine assembles coefficients through a
$\Gamma$-point scalar and emits no kernel block at all, and by
Theorem~\ref{thm:obstruction} $\Gamma$-point data are structurally incapable of
constraining $C_{\mathrm{shp}}$. This is a statement about the instrument, made
by reading it, not a report of a failed run.
\chk{FINDING: the marked-cluster engine emits the Gamma scalar only — a completed 609-sweep cannot decide C\_shp}

\subsection{The tier collapse, as two integers}

The two shape coefficients $B$ and $D$ vanish in the recorded kernel even
though the two-hop enumeration generates both, and the degree count does
\emph{not} forbid them. The carrier numerator reaches $a$-degree three, because
the projection contributes one further power of $a$: for any vector whose
components have squared modulus $a_i$ --- the carrier is one, in whichever of
the two equivalent labellings --- the quadratic form $\mathrm{diag}(a_i^2)$
evaluates to $\sum_i a_i^3=q^3-3qe_2+3e_3$, an explicit witness carrying a
nonzero $e_3$. An earlier degree-bound mechanism proposed
for the collapse was therefore retracted. What is checked is finer: every
degree-three block amplitude is an integer multiple of one raw record weight
$x$, and
\[
  B=0:\;(+1,+1,-2)\cdot x,\qquad
  D=0:\;(-3,+6,-3)\cdot x
\]
over the displacement shells. The kernel does populate the forbidden tier; the
dynamics cancels it with small integers. Any mechanism must reproduce those two
vectors.
\chk{RETRACTED: the vertex count does NOT forbid B\_shp and D\_shp}
\chk{FINDING: the tier collapse is two integer cancellations at record level}
\chk{the vanishing coefficients are exactly the degree-3 ones}

\subsection{A near-$\Gamma$ statement that does not adjudicate}

One quantitative fourth-order statement can be made without adjudicating
anything, because it depends on the kernel only through $\sqrt{W_4}$. What the
criterion needs is a lower bound on $q$ over the region $|k|\ge r$, and there it
is sharp: writing $q=\sum_mh(k_m^2)$ with $h(t)=2(1-\cos\sqrt t)$, the identity
$(s\cos s-\sin s)'=-s\sin s\le0$ makes $h$ concave and increasing on
$[0,\pi^2]$, so the minimum over the simplex sits at a vertex and
\begin{equation}\label{eq:shellmin}
  \min_{|k|\ge r}q(k)=4\sin^2(r/2),
\end{equation}
attained on a coordinate \emph{axis}. With $\tau(u)\ge t_3u^2$ for $u>0$, the
fourth-order shape terms are therefore dominated by the second-order dispersion
outside a ball of radius $2\arcsin\bigl(\tfrac u2\sqrt{W_4/(\theta t_3)}\bigr)$,
which is $Ku$ to leading order with $K=\sqrt{W_4/(\theta t_3)}$: the two records
give $K=10.85$ and $K=15.06$ --- floating point, since one of the two bandwidths
exists only as a float --- and the larger bounds both, so the exclusion radius
can be stated for either record without choosing between them. Jordan's
inequality $q(k)\ge(4/\pi^2)|k|^2$ also holds on the whole zone and is tight at
$k=0$ and at the corner, but the corner is not where this bound is used; on the
axis it is loose by $\pi^2/4$, and the constants it gives, $17.04$ and $23.66$,
are exactly $\pi/2$ times the sharp ones. It is not a bound on the open
coefficient: a value of $C_{\mathrm{shp}}$ outside the two records carries its
own $W_4$, and nothing here confines the answer to the interval the records
span. The statement is non-vacuous only for $u<0.133$ at $\theta=1/2$, and that
threshold is the same for both constants --- the sharp radius reaches $\pi$ when
$\tfrac u2\sqrt{W_4/(\theta t_3)}=1$, which is exactly where $Ku=\pi$ for the
Jordan constant, because the two bounds agree at the zone corner. The excluded
zone fraction grows as $u^3$.
\chk{Jordan bound q(k) >= (4/pi\^{}2)|k|\^{}2 holds on the whole zone}
\chk{the exact zone minimum of q on |k| >= r is 4 sin\^{}2(r/2), and Jordan overstates it by pi/2}
\chk{the criterion survives C2: K depends on the kernel only through sqrt(W\_4)}
\chk{the statement is non-vacuous only below an explicit coupling}

\section{External comparisons}
\label{sec:external}

External results count here because they are \emph{independent} --- produced
without knowledge of this program --- not because they are published. Each is
an assertion until something checks it.

\paragraph{Hamer 1989, both sectors.} The $1^{+-}$ rest-frame series agrees
with $E_{\mathrm{flat}}$ at every shared order: $16/3\div2=8/3$ exactly,
$a_1=+1$, and $2a_2$, $4a_3$ matching $11/306$ and $d_3$ to $9.9\times10^{-14}$
and $1.1\times10^{-12}$. The $0^{++}$ series matches the charge-even $A_1^{++}$ $\Gamma$-point
coefficients at $x^2$ and $x^3$, with the sign structure ($a_1=-1$ in the
scalar channel where the carrier's is $+1$) reproduced; the $E^{++}$ doublet is
not constrained by it, which is why Proposition~\ref{prop:evengamma} needs the
splitting and not the level. The conversion $m_n=2^{n-1}a_n$ is $x=2u$ together
with Hamer's normalisation $ma=(g^2/2)M$ --- the $2^n$ from the substitution,
the half from the normalisation --- exactly, with no $4^r$ ambiguity anywhere.
Hamer states in print that his $x^3$ and $x^4$ terms disagree with the earlier
series of~\cite{KSS1976}; that paper has not been read here, so the
disagreement is recorded and not adjudicated.
\chk{Hamer's 1+- series matches the C-odd Gamma-point coefficients through x\^{}3}
\chk{Hamer's 0++ series matches the C-even Gamma-point coefficients through x\^{}3}
\chk{the m\_n = 2\^{}(n-1) a\_n bridge is the x = 2u conversion}
\chk{the Hamer table is pinned, and the a\_4 agreement is primary-source}

\paragraph{The string tension.} The Kogut--Pearson--Shigemitsu $SU(3)$
table~\cite{KPS1981} --- read here from the 1980 preprint, the journal pages
themselves being unread --- equals this program's certified $\sigma$ series
exactly,
$2^{n-1}t_n=\sigma_n$ for $n=2,\dots,5$, in the physical sign convention
$\sigma_n^{\mathrm{phys}}=(-1)^n\sigma_n^{\mathrm{raw}}$. The same series
reproduces $E_{\mathrm{flat}}$ term by term through $u^3$, including the
discriminating $u^3$ coefficient.
\chk{the KPS 1980 string-tension table equals the certified sigma series EXACTLY}
\chk{sigma\_n\^{}phys = (-1)\^{}n sigma\_n\^{}raw, and C5 is the n = 3 case}
\chk{the ratio and sigma series reproduce E\_flat exactly}

\paragraph{Euclidean series.} Two printings of the same M\"unster/Seo series
--- the second erratum to~\cite{Munster1981}, which prints the definitive
corrected table and credits Ukawa and Seo~\cite{Seo1982} with locating the
error, and Smit's Table~1 scalar column~\cite{Smit1982}, which credits them ---
agree rational for rational through $u^8$. That is agreement
transcription for transcription, not two independent computations, and it is
recorded as such. One shift is recorded and not adjudicated: M\"unster's 1985
table~\cite{Munster1985} differs from the 1982 erratum row at eighth order by
exactly $-96$, a fourth correction already in print by 1983~\cite{Munster1983}
that no erratum records. This is a Euclidean record; no Hamiltonian comparison is drawn from
it, and none is permitted.
\chk{the errata-resolved Euclidean series is doubly sourced, transcription for transcription}
\chk{FINDING: Munster's 1985 table shifts his 1982 erratum at eighth order}

\paragraph{The overlap obstruction.} Schierholz~\cite{Schierholz1988} measured
in 1988 that at $\beta=5.9$ --- in the Euclidean scaling region, not at strong
coupling --- a bare local operator carries only a couple of percent of the state
it interpolates, and gave the scaling $a^5$ for a local dimension-four operator.
Nothing is transferred from that measurement to here: it is a Euclidean Monte
Carlo result at the opposite end of the coupling axis from $u\ll1$, and the
firewall of Section~\ref{sec:scope} forbids importing it as a bound. Read as a
warning rather than as a bound, it is why \texttt{G18} must be posed in a
smeared basis: an overlap that falls as a power of the spacing is not something
better statistics recover.
\chk{the overlap obstruction was published in 1988, and it scales}

\paragraph{Cross-regime discipline.} Two indexed papers declare an explicit
scope firewall and neither supplies a value here:
Davoudi--Raychowdhury--Shaw~\cite{DRS2021}, whose Hilbert-space cost comparison
is locked to $1{+}1$ dimensions with matter, and
Kadam--Raychowdhury--Stryker~\cite{KRS2023}. The rule is enforced by the
literature validator rather than recorded in prose.
\chk{a cross-regime paper never supplies a value}

\section{Scope}
\label{sec:scope}

This work does not establish a four-dimensional infinite-volume Yang--Mills
mass gap, a continuum glueball mass, convergence of the strong-coupling series
to the continuum, regulator independence of the first dispersive order, the
full $SU(3)$ fourth-order rest coefficient or off-axis shape, or an
identification of any level here with a physical glueball. At $\Gamma$ the
three face orbitals transform as an axial vector of the cubic group, which
with parity even and charge conjugation odd gives the retained-space label
$T_1^{+-}$; this is an operator-space assignment, not a spectral one.

The carrier is macroscopically degenerate at finite $L$, but its nearest
incidence separation is $\Delta_L(u)=4\tau(u)\sin^2(\pi/L)+O(u^4)$, which
falls as $L^{-2}$ and so provides no volume-uniform spectral isolation.
\chk{Delta\_L = 4 tau(u) sin\^{}2(pi/L) is positive and falls as L\^{}-2}

Three crossings are named explicitly, and none of them is made here.

\begin{itemize}\itemsep2pt
\item \textbf{Finite lattice to infinite volume} requires an inhomogeneous
  Wilson free-energy bound with a useful $\log K_\alpha$ and a source-radius
  reduction. Both are open, and are carried in the accompanying gap register as
  \texttt{G17}.
\item \textbf{Protected level to particle} requires a volume-uniform overlap
  of a smeared, dressed $T_1^{+-}$ operator built on the carrier with the
  transfer-matrix spectrum, and multi-plaquette survival of the protected
  level (\texttt{G18}). The published overlap measurement of
  Section~\ref{sec:external} is why this theorem must live in the smeared
  basis.
\item \textbf{Lattice to continuum} (\texttt{G19}) is untouched. Small $u$ is
  strong coupling; the continuum is the opposite end of the axis.
\end{itemize}

\noindent The second and third are the only route from the statements above to
any sentence containing the words ``mass gap'', and nothing in this edition
shortens that route.

Four premises earlier in the paper are arguments in prose rather than results,
and all four are load-bearing. None has a check.

That the range of $P_-$ is exactly the charge-odd one-plaquette span for
$L\ge3$ is argued in Section~\ref{sec:sector} from the link count of a reduced
fundamental loop.

That one-shared-link processes exhaust the second-order \emph{off-diagonal}
intermediates is assumed in Theorem~\ref{thm:allrank}. Three neighbouring
statements \emph{are} settled and should not be confused with it: the channel
list is complete (Lemma~\ref{lem:channels}), two distinct faces of this lattice
share at most one link, and the vacuum-mediated route vanishes in this sector
(Proposition~\ref{prop:vacuum}). The gap is in the process enumeration, not in
the channel list, and Theorem~\ref{thm:weingarten} is silent about it.

That the degree-$(2,2)$ Haar moment over $SU(N)$ equals the $U(N)$ one is
argued in the proof of Theorem~\ref{thm:weingarten} from the $U(1)$ invariance
of the integrand and the exhaustion of the invariants below the rank. The
argument is standard and the range is the theorem's own, but nothing here
checks it.

That the charge-even sector's second-order Bloch symbol is the unsigned
incidence $\mathcal N(k)$ --- the signed incidence with the boundary
orientations dropped --- is asserted in Section~\ref{sec:second} and used
throughout Section~\ref{sec:even} and Appendix~\ref{app:even}. The registry
records it as a convention rather than a result, and the closed-form cubic, the
range, the bandwidth and the $\Gamma$-splitting all rest on it.

\section{Reproducibility}
\label{sec:repro}

A portable standard-library Python program, \texttt{paper/verify\_core.py} in
the accompanying repository (the root-level file of the same name belongs to
the earlier flat-band manuscript and is a different, smaller program), runs
twenty-five checks in nine claim groups: the rational coefficient ledger; the
Weingarten derivation of Theorem~\ref{thm:weingarten} by explicit index
summation; the incidence identity, \emph{decided} in exact Gaussian-rational
arithmetic at rational torus points rather than checked to a residual; the
chain condition $\partial_2\partial_3=0$ over $\mathbb Z$ and the finite-torus
ranks with an explicit kernel basis, together with the closed cube surface of
Proposition~\ref{prop:cube}; the Rayleigh quotient of the wrapping sheets; the
two-cube connected geometry rebuilt from oriented cell boundaries, its exact
$B=6$ and $B=4$ spectra, the endpoint Casimir budgets and the weak retention
rule they decide, and Reported result~\ref{rep:census} channel by channel; the charge-even cubic of Theorem~\ref{thm:evencubic} with $p+q=12$ and
the four high-symmetry spectra of both sectors and the $\Gamma$ splitting of
Proposition~\ref{prop:evengamma}; the checkpoint algebra of
\eqref{eq:checkpoints}, the second witness, and the vanishing of $e_2$ on
every axial cut;
and the coupling convention, including the $16$ and $64$ by which a $4^r$
rescaling misses. Every value in it is a
\texttt{Fraction}: no float is constructed anywhere, so ``agreement'' there
always means equality. It has no dependencies and runs in about half a second.

Nothing in \texttt{verify\_core.py} assumes or decides a value of
$C_{\mathrm{shp}}$. Its fourth-order group is the obstruction as arithmetic:
the extraction formulas, the blindness of $X$, and the identical vanishing of
$e_2$ on axial cuts.

Every theorem, proposition, corollary and reported result above names the
machine check that establishes it, and the appendix tables name checks too.
Three displayed objects deliberately carry none, because none is a result of
this program: the dimension and Casimir table \eqref{eq:table}, which is
standard $SU(N)$ representation theory; the boundary formulas of
Appendix~\ref{app:chain}, which are definitions --- what is checked there is
that they compose to zero; and Table~\ref{tab:counters}, which is provenance
about a commit rather than a claim about the physics, and whose rows are
re-measured by the commands printed beside it. The checks live in the
accompanying repository, where
\texttt{workhouse verify --only '<name>'} re-runs any one of them alone, with
its numbers, in about a second. That device is itself checked: a registered
invariant reads every pinned edition in \texttt{paper/} --- this one included
--- resolves every printed label against the live registry, and fails if one is
renamed, deleted, or left red. The guard is the reason this edition's labels
can be trusted at all: the same invariant read only the previous edition until
this one landed, and a label that had drifted passed unnoticed until it was
made to read both.

This edition pins commit
\texttt{af3ab5f2e9860e75372c1a143a33c828f8b0d9ba} of the accompanying
repository --- the commit at which every check name printed above resolves,
and the last one before this file itself lands. It carries the counters in
Table~\ref{tab:counters}. They describe \emph{that} commit and no other, and
each row has its own command: after
\texttt{git checkout af3ab5f}, \texttt{make verify} gives the $T_1$ and $T_2$
rows and \texttt{make check} the test count, while the $T_0$ row is the
declaration scrape that \texttt{workhouse certified} prints ---
\texttt{make lean} compiles the tree and reports no count.
They are provenance rather than independent evidence; the argument is above, and the checks are the checks.

\begin{table}[t]
\centering
\begin{tabular}{@{}lr@{}}
\toprule
Layer & Count \\
\midrule
$T_0$ — Lean 4 theorems, no \texttt{sorry} & 40 \\
$T_1$ — exact re-derivations & 188 \\
$T_2$ — numerical, within a stated tolerance & 49 \\
Repository tests & 486 \\
\bottomrule
\end{tabular}
\caption{Counters at the pinned commit. The $T_0$ row counts theorem
declarations with no \texttt{sorry}, which is a scrape of the tree;
\texttt{make lean} is what compiles them, and the standard the tree is held to
is that every axiom footprint lies in $\{$\texttt{propext},
\texttt{Classical.choice}, \texttt{Quot.sound}$\}$. That compilation was last
run elsewhere: the environment this edition was typeset in carries no Lean
toolchain, and the row is not asserted as re-measured here.}
\label{tab:counters}
\end{table}

A smaller layer is proof-checked rather than merely re-derived. The Lean
development compiles the rank law with denominators cleared, $t_3=5/612$ and
$t_2=0$, the third-order ledger identity, the axial law at every exceptional
rank, the cycle count $(L^3-1)+3=L^3+2$, the four checkpoint-extraction
formulas and the deltas they invert, and four statements about the charge-even
cubic including the identity that its middle coefficient is $p^2$. It is
exact-rational and polynomial algebra, plus the trigonometric checkpoint
evaluations over $\mathbb R$ that supply those deltas --- the $\mathbb R$
section exists because the extraction theorems formalised only half a statement
without it. There is no perturbative derivation and no operator theory.

An earlier edition added that ``nothing that bears on $C_{\mathrm{shp}}$
appears there''. That was false, and false about the one genuinely open item
here. The Lean core defines the historical side
$C_{\mathrm{old}}=-211835444920651/4405310420659200$ as an exact rational and
proves three internal-consistency relations about it: $C$ from $\beta$, the
width as $\alpha+\beta$, and the tier-collapse relation. What it contains no
spelling of, anywhere in the tree, is the v10a.26 value --- so the layer
touches $C_{\mathrm{shp}}$ and still prefers neither side, and that
even-handedness is an absence checked by inspection rather than asserted.
\chk{the Lean core DOES carry the historical C\_shp side, and still adjudicates nothing}

Machine checks certify algebra relative to stated definitions and inputs; they
do not replace a derivation, and they do not certify an uncompleted off-axis
contraction --- which is why Section~\ref{sec:fourth} stops where it does.

\appendix

\section{Channel weights in closed form}\label{app:closed}

Substituting \eqref{eq:table} into \eqref{eq:w},
\begingroup\small
\begin{align*}
  w_{\mathbf 1}&=-\frac{2}{N(N^2-1)}, &
  w_{\Lambda^2F}&=-\frac{N-1}{(N+1)(2N-3)},\\
  w_{\mathrm{Adj}}&=-\frac{2(N^2-1)}{N(2N^2-1)}, &
  w_{\mathrm{Sym}^2F}&=-\frac{N+1}{(N-1)(2N+3)} .
\end{align*}
\endgroup
Their family sums are \eqref{eq:AN} and \eqref{eq:BN}, and the difference has
numerator $2N(N^2-4)$, which is the $SU(2)$ zero and the positivity for
$N\ge3$ at once.
\chk{the four channel weights follow from dimension and Casimir} At $N=3$:
$w_{\mathbf 1}=-1/12$, $w_{\mathrm{Adj}}=-16/51$, $w_{\Lambda^2F}=-1/6$,
$w_{\mathrm{Sym}^2F}=-2/9$, which are the four values
Reported result~\ref{rep:census} matches against \eqref{eq:census}.

\section{The Weingarten index sums}\label{app:wg}

With $\mathrm{Wg}(e)=1/(N^2-1)$ and $\mathrm{Wg}((12))=-1/(N(N^2-1))$, and
$\sigma,\tau$ running over $S_2$,
\[
  \begin{aligned}
    &\bigl\langle U_{i_1j_1}U_{i_2j_2}
       \bar U_{i'_1j'_1}\bar U_{i'_2j'_2}\bigr\rangle\\
    &\qquad=\sum_{\sigma,\tau}\mathrm{Wg}(\sigma\tau^{-1})
      \prod_a \delta_{i_ai'_{\sigma(a)}}\delta_{j_aj'_{\tau(a)}} .
  \end{aligned}
\]
Summing the delta products over all $N^4$ index quadruples,
\[
  M_{\mathrm{direct}}=\mathrm{Wg}(e)(N^4+N^2)+2\,\mathrm{Wg}((12))N^3=N^2,
\]
\[
  M_{\mathrm{cross}}=2\,\mathrm{Wg}(e)N^3+\mathrm{Wg}((12))(N^4+N^2)=N .
\]
Both are confirmed by explicit summation at $N=3,4,5$.
\chk{the shared-link weights are Weingarten, not an isotropy assumption}
The related fourth moment $\int|U_{11}|^4\,dU=2/[N(N+1)]$ is $1/6$ at $N=3$ and in particular
nonzero, which is what refuted a claimed Haar vanishing in the balanced
$(2,2)$ sector.
\chk{SU(3) Weingarten values follow from the general formula}
\chk{the fourth moment integral |U\_11|\^{}4 = 1/6 at N = 3}

\section{The $SU(3)$ coefficient ledger}\label{app:ledger}

\begin{center}
\begin{tabular}{@{}lll@{}}
\toprule
Order & Quantity & Value\\
\midrule
$u^0$ & electric plaquette & $8/3$\\
$u^1$ & scalar & $1$\\
$u^2$ & $BB\dg$ coefficient & $5/612$\\
$u^2$ & flat scalar & $11/306$\\
$u^3$ & $BB\dg$ coefficient & $1975/124848$\\
$u^3$ & per-neighbour leakage & $-12331/249696$\\
$u^3$ & flat scalar & $-109151/249696$\\
\bottomrule
\end{tabular}
\end{center}
\chk{the manuscript's SU(3) ledger is this registry, value by value}

The $u^1$ row is $SU(3)$-only rather than an $SU(3)$ specialisation: it is
$-\langle p,-|V|p,-\rangle=1$, which needs the $\epsilon$ channel and vanishes
for $N\ge4$. Of the rest, the $u^0$ row and the $u^2$ $BB\dg$ coefficient
specialise all-rank statements ($2\CF$ and $t_N$); the $u^2$ flat scalar
inherits the $SU(3)$ tower term through \eqref{eq:assembly}; and the three
$u^3$ rows are $SU(3)$-only and conditional on the supplied domino ledger,
exactly as Section~\ref{sec:su3} says.

\section{The charge-even ledger}\label{app:even}

In the same convention, with $\lambda\in\{12,-4\}$ and \eqref{eq:assembly}. The
$\lambda=-4$ rows and the bandwidth read off them are even-$L$,
thermodynamic-limit values: $\lambda=-4$ requires some $k_j=\pi$, which the
periodic grid samples only at even $L$. The $\lambda=12$ rows sit at $\Gamma$
and hold on every grid.

\begin{center}
\begin{tabular}{@{}lll@{}}
\toprule
Order & Quantity & Value\\
\midrule
$u^2$ & hopping $t_{2,+}=\ell_3$ & $-11/306$\\
$u^2$ & on-site diagonal & $223/1020$\\
$u^2$ & band bottom ($\lambda=12$) & $-217/1020$\\
$u^2$ & band top ($\lambda=-4$) & $1109/3060$\\
$u^2$ & bandwidth $16|t_+|$ & $88/153$\\
$u^3$ & hopping $t_{3,+}$ & $-6335/249696$\\
$u^3$ & on-site diagonal ($\lambda=0$) & $52163/260100$\\
$u^3$ & band bottom ($\lambda=12$) & $-54049/520200$\\
$u^3$ & band top ($\lambda=-4$) & $471353/1560600$\\
\bottomrule
\end{tabular}
\end{center}

\chk{one assembly formula gives every C-even value at both orders}

Because $t_{r,+}<0$ at both orders, the $\lambda=-4$ endpoint is the band
\emph{maximum} in this sector. A certificate key naming it \texttt{bandmin} at
both orders is a naming error against its own arithmetic, recorded rather than
silently corrected.
\chk{FINDING: the certificate key 'bandmin' holds the band MAXIMUM, at both orders}

\section{Finite-volume chain calculation}\label{app:chain}

For $i<j$,
\[
  \partial_2[x;i,j]=[x+e_i;j]-[x;j]-[x+e_j;i]+[x;i],
\]
\begin{align*}
  \partial_3[x;1,2,3]=\;&[x+e_1;2,3]-[x;2,3]\\
  &-[x+e_2;1,3]+[x;1,3]\\
  &+[x+e_3;1,2]-[x;1,2].
\end{align*}
Substituting the second into the first cancels every edge pairwise, so
$\partial_2\partial_3=0$ over $\mathbb Z$.
\chk{d\_2 d\_3 = 0 on the built complex}
On the torus $\ker\partial_3$ is generated by the sum of all oriented cubes,
so $\operatorname{rank}\partial_3=L^3-1$.
\chk{rank d\_3 = L\^{}3 - 1 on the built complex}

\section{The two-cube geometry}\label{app:twocube}

The eleven oriented plaquettes of the $(3,2,2)$ prism, by base vertex and
ordered plane axes, are
\[
\begin{aligned}
&(000;12),\;(000;13),\;(000;23),\;(001;12),\\
&(010;13),\;(100;12),\;(100;13),\;(100;23),\\
&(101;12),\;(110;13),\;(200;23),
\end{aligned}
\]
with left, right and shared-face index sets
$I_L=(0,1,2,3,4,7)$, $I_R=(5,6,7,8,9,10)$, $I_F=(7)$. Building $G=BB^{\mathsf
T}$ from the oriented boundaries and applying the geometric M\"obius fold
\eqref{eq:mobius} leaves exactly four nonzero off-diagonal pairs,
$(0,5),(1,6),(3,8),(4,9)$, each equal to $-1$: four disjoint $2\times2$ blocks
and three isolated directions, which is why the spectra in
Section~\ref{sec:twocube} are closed form once the diagonals below are
supplied. The two connected diagonals are
\[
  D_{B6}=\tfrac1{612}\operatorname{diag}
  (-2317,-2317,-2295,-2317,-2317,
\]
\[
  \qquad -2317,-2317,0,-2317,-2317,-2295),
\]
\[
  D_{B4}=\operatorname{diag}
  (-\tfrac74,-\tfrac74,-\tfrac{15}4,-\tfrac74,-\tfrac74,
\]
\[
  \qquad -\tfrac74,-\tfrac74,0,-\tfrac74,-\tfrac74,-\tfrac{15}4),
\]
so both spectra are one line of arithmetic: a pair block contributes
$D_{ii}\mp c$ twice and each isolated direction its own $D_{ii}$. The $B=4$
comparator on the same geometry gives $\{0,(-\tfrac{15}4)^{2},(-\tfrac{11}6)^{4},(-\tfrac53)^{4}\}$:
the paired levels sit \emph{above} the isolated $-15/4$ pair, where at $B=6$
they sit below it. The three isolated levels $\{-\tfrac{15}4,-\tfrac{15}4,0\}$
are identical in the two truncations, so the entire effect of the cutoff lives
in the four cross-cell blocks. That reversal is not the coefficient's sign by itself --- within a
$2\times2$ block the pair splits as $D\mp c$, which is invariant under
$c\to-c$ --- but the truncation's effect on the connected diagonal as well.
What the sign of $c$ alone does reverse is the ordering on a \emph{single}
cube, where the kernel is $\alpha I+tG$ with one diagonal and the spectrum of
$G$ fixed; that is the shell reversal the one-cube reconstructions report.
\chk{the B=4 comparator on the same geometry gives -1/12 with the reversed spectrum}

\begin{thebibliography}{99}\small

\bibitem{KS1975} J.~B. Kogut and L.~Susskind,
  \emph{Hamiltonian formulation of Wilson's lattice gauge theories},
  Phys. Rev. D \textbf{11}, 395 (1975).

\bibitem{KSS1976} J.~B. Kogut, D.~K. Sinclair and L.~Susskind,
  \emph{A quantitative approach to low-energy quantum chromodynamics},
  Nucl. Phys. \textbf{B114}, 199 (1976).

\bibitem{Weingarten1978} D.~Weingarten,
  \emph{Asymptotic behavior of group integrals in the limit of infinite rank},
  J. Math. Phys. \textbf{19}, 999 (1978).

\bibitem{KPS1981} J.~B. Kogut, R.~B. Pearson and J.~Shigemitsu,
  \emph{The string tension, confinement and roughening in $SU(3)$ Hamiltonian
  lattice gauge theory}, Phys. Lett. B \textbf{98}, 63 (1981).

\bibitem{Munster1981} G.~M\"unster,
  \emph{Strong coupling expansions for the mass gap in lattice gauge theories},
  Nucl. Phys. \textbf{B190}, 439 (1981); Erratum and Addendum
  \textbf{B200}, 536 (1982); Erratum \textbf{B205}, 648 (1982). The two errata
  were read here; the main paper was not.

\bibitem{Seo1982} K.~Seo,
  \emph{Glueball mass estimate by strong coupling expansion in lattice gauge
  theories}, Nucl. Phys. \textbf{B209}, 200 (1982). Not read here; cited as the
  independent locator of the eighth-order error.

\bibitem{Smit1982} J.~Smit,
  \emph{Estimate of glueball masses from their strong coupling series in
  lattice QCD}, Nucl. Phys. \textbf{B206}, 309 (1982).

\bibitem{Munster1983} G.~M\"unster,
  \emph{Physical strong coupling expansion parameters and glueball mass
  ratios}, Phys. Lett. B \textbf{121}, 53 (1983). Not read here; cited as the
  1983 printing the 1985 shift is consistent with.

\bibitem{Munster1985} G.~M\"unster,
  \emph{Effective transfer matrix for low-lying glueball states in lattice
  gauge theory}, Nucl. Phys. \textbf{B256}, 67 (1985).

\bibitem{Schierholz1988} G.~Schierholz,
  \emph{Glueball masses in $SU(3)$ lattice gauge theory},
  Nucl. Phys. B (Proc. Suppl.) \textbf{4}, 11 (1988).

\bibitem{Hamer1989} C.~J. Hamer,
  \emph{Hamiltonian strong coupling expansions for glueball masses in
  $SU(3)$}, Phys. Lett. B \textbf{224}, 339 (1989).

\bibitem{CM2003} J.~Carlsson and B.~H.~J. McKellar,
  \emph{$SU(N)$ lattice gauge theory on a single cube}, hep-lat/0303022 (2003).

\bibitem{CollinsSniady2006} B.~Collins and P.~\'Sniady,
  \emph{Integration with respect to the Haar measure on unitary, orthogonal and
  symplectic group}, Comm. Math. Phys. \textbf{264}, 773 (2006),
  math-ph/0402073.

\bibitem{KRS2023} S.~V. Kadam, I.~Raychowdhury and J.~R. Stryker,
  \emph{Loop-string-hadron formulation of an $SU(3)$ gauge theory with dynamical
  quarks}, Phys. Rev. D \textbf{107}, 094513 (2023), arXiv:2212.04490.

\bibitem{CBB2026} A.~N. Ciavarella, I.~M. Burbano and C.~W. Bauer,
  \emph{Efficient truncations of $SU(N_c)$ lattice gauge theory for quantum
  simulation}, arXiv:2503.11888 (2025). No journal reference appears on the
  arXiv record as of 30 August 2026.

\bibitem{Balaji2026} S.~Balaji \emph{et al.},
  \emph{Quantum simulation of $SU(3)$ lattice Yang--Mills theory: one-cube
  benchmarks and circuit resource estimates}, arXiv:2509.25865. Used here for
  its public \texttt{ymcirc} $d=3$, $B=4$ plaquette table; the paper itself has
  not been read.

\bibitem{DRS2021} Z.~Davoudi, I.~Raychowdhury and A.~Shaw,
  \emph{Search for efficient formulations for Hamiltonian simulation of
  non-Abelian lattice gauge theories}, Phys. Rev. D \textbf{104}, 074505
  (2021).

\bibitem{Sutherland1986} B.~Sutherland,
  \emph{Localization of electronic wave functions due to local topology},
  Phys. Rev. B \textbf{34}, 5208 (1986). Not read here.

\bibitem{Mielke1991} A.~Mielke,
  \emph{Ferromagnetism in the Hubbard model on line graphs and further
  considerations}, J. Phys. A \textbf{24}, 3311 (1991). Not read here.

\bibitem{BWB2008} D.~L. Bergman, C.~Wu and L.~Balents,
  \emph{Band touching from real-space topology in frustrated hopping models},
  Phys. Rev. B \textbf{78}, 125104 (2008), arXiv:0803.3628. Not read here.

\bibitem{Hazra2026} T.~Hazra,
  \emph{New class of exactly flat topological bands: compact localised states
  protected by local graph topology}, arXiv:2607.22831.

\end{thebibliography}

\end{document}
